% !TEX root = ../main.tex
\chapter{Implementation and Empirical Results}
\label{ch:results}

Chapter~\ref{ch:methodology} described how EndorseRank and the Adaptive Weighted PageRank (AWP) baseline are built, which construct checks are applied, how uncertainty is quantified and how runtime is measured. This chapter runs that design on Arbitrum One and reports what was observed. It is a single case study. There is one chain and one observation window, with one matched cohort for every same-window comparison, a spender cohort and a trader cohort for the temporal holdouts, and one expanded address pool for node-count sensitivity. A registered replication on labels from June to August~2026 closes Section~\ref{sec:holdout-results}. Every table is generated from the evaluation summaries of Section~\ref{sec:reproducibility}, and each Kendall $\tau$ coefficient is reported with its confidence interval.

\dissertationSubheading[sec:dataset]{Case study setting: Arbitrum One}

Arbitrum One was chosen as the case because it is a heavily used Ethereum rollup with mature DeFi protocols (Section~\ref{sec:ethereum-l2}), because its logs are included in the public BigQuery export that the pipeline reads, and because one of its margin-trading protocols, GMX~V2, supplies the wallet seed set described below. The observation window covers 1~December~2025 to 31~May~2026. Six months is long enough to span several market regimes and to let the tenure and active-month proxies vary. It is also short enough that the time weight of an event never falls below about half of the weight of a recent one (Figure~\ref{fig:logistic-decay-methods}).

\dissertationSubsubheading[sub:cohorts]{Cohorts}

Three address sets are used, each with its own purpose.

\begin{itemize}
\item Matched cohort ($n=5{,}521$). These are the GMX~V2 margin-trading wallets on Arbitrum One with at least three closes (\texttt{PositionDecrease} events, partial closes and liquidations included) between 1~December~2025 and 31~May~2026. The list is fixed from GMX data before any approval or transfer is extracted, and the BigQuery approval and transfer queries are filtered to it. Each method's evaluation graph, $G_{\text{approve}}$ for the latest allowances and $G_{\text{transfer}}$ for the transfers over the window, keeps every edge with at least one endpoint in the cohort. Membership does not depend on either graph. No allowance edge touches $966$ cohort wallets, no transfer edge touches $381$, and $370$ are touched by neither; a wallet that no edge of a graph touches is absent from it and scores zero under that method. Every result in the same-window alignment, difference-test and benchmark sections is computed on this cohort.
\item Holdout spender cohort ($n=1{,}335$). It contains every spender that holds at least one positive latest allowance at $t_1$, the end of the scoring period (28~February~2026, 23:59:59~UTC), in the extracted approval data. Most of these spenders are contracts. When their code was read from an Arbitrum node at the end of September 2026 (block 510{,}408{,}687), 1,174 of the 1,335 addresses held contract code, and the other 161 were externally owned accounts, one of which delegates to contract code under EIP-7702. All 100 spenders that EndorseRank ranks highest are contracts, and so are 221 of the 222 that gained a new approval pair from March to May. A BigQuery check of the scoring window agrees where it can decide: every spender that emitted an event log holds code. It cannot decide for the 686 spenders that neither emitted a log nor sent a transaction in the window. Of these, 530 are contracts, among them routers whose events are emitted by other contracts, such as the GMX exchange router named in the pipeline configuration, and 156 are externally owned accounts. Scores for these wallets are tested against owner--spender approval pairs and transfer senders first seen in the holdout period, 1~March~2026 to 31~May~2026 (Section~\ref{sec:holdout-protocol}).
\item Expanded wallet pool ($N=99{,}080$). It combines the matched cohort with $93{,}559$ further addresses sampled from the chain's monthly ERC-20 activity. The pool is used for two analyses only, the node-count sensitivity of runtime and the sensitivity to sample definition (Section~\ref{sub:expanded-pool}). The added addresses bring no edges beyond those extracted for the cohort, so the full pool's solver graph is the cohort graph, with the same nodes and edges; the relevant results section restates this caveat.
\end{itemize}

The matched-cohort graph sizes are listed with the benchmark figures in Table~\ref{tab:benchmark-runtime}. On the same wallets, the endorsement graph has roughly one-ninth as many edges as the transfer graph. Both graphs are drawn from the same wallets and the same window, so the imbalance lies in what each graph records; it is not a sign of an error in construction. EndorseRank keeps only the latest non-zero allowance for each token, owner and spender and adds one weighted term per token, so repeated approvals collapse into one directed edge per owner--spender pair. AWP adds a weighted term for every transfer between two wallets. Whatever differences in score and alignment separate the methods trace back to their edges and restarts, since the sampling is the same for both.

\dissertationSubsubheading[sub:implementation]{Implementation}

Every algorithm runs on one Python implementation of PageRank, with $d=0.85$, tolerance $\varepsilon=10^{-8}$ and an iteration cap of $I_{\max}=300$. EndorseRank and AWP differ in their edge lists and in the restart vector, which is uniform for EndorseRank and follows sending activity for AWP. The coupled operator mixes its two edge lists node by node, as in Equation~\eqref{eq:coupled-transition} (Section~\ref{sec:implementation-notes}). Each run writes its scores, per-proxy statistics, bootstrap confidence intervals and timings to one file. The tables of this chapter come from runs made on 2~October~2026, except those of the registered replication, which was evaluated on 1~October; Appendix~\ref{ch:appendix-eval} lists the runs, the repository revision, the configuration file and the summary files.

\dissertationSubheading[sec:benchmark-results]{Computational benchmark (Claim~C1)}

Claim~C1 is that EndorseRank costs less to compute than AWP on the same input. Following the timing procedure of Section~\ref{sec:computational-benchmark}, Table~\ref{tab:benchmark-runtime} reports for each method, on the matched cohort, the PageRank solve time and its standard deviation over the five timed repeats, the peak memory of the warm-up solve, the number of iterations to convergence and the size of the evaluation subgraph.

% Real BigQuery data -- regenerate with:
%   A-Skill-Programs/margin_rank/scripts/run_dissertation_eval.py --real --export-latex
% Generated by export_latex_results.py -- do not edit by hand unless re-export fails
\begin{table}[htbp]
\centering
\caption[Computational benchmark on the matched wallet cohort.]{Computational benchmark on the matched wallet cohort ($n=5{,}521$; AWP/EndorseRank runtime ratio 10.2$\times$). Runtime = mean wall-clock seconds over 5 timed PageRank solves after one untimed warm-up; SD over the same 5 runs. $|V|$, $|E|$ = nodes and edges of each method's evaluation subgraph. C-PR walks on the union of both layers, so $|E|$ is the sum of the two edge sets; S-PR runs on the transfer layer and its runtime includes the EndorseRank solve that supplies its restart vector.}
\label{tab:benchmark-runtime}
\fitwidth{%
\begin{tabular}{lrrrrrr}
\toprule
Method & Runtime (s) & SD (s) & Peak mem.\ (MB) & Iters & $|V|$ & $|E|$ \\
\midrule
EndorseRank & 0.401 & 0.016 & 4.6 & 37.0 & 6{,}442 & 14{,}727 \\
AWP & 4.105 & 0.193 & 38.2 & 100.0 & 30{,}791 & 137{,}087 \\
\midrule
C-PR ($\lambda=0.5$) & 4.760 & 0.471 & 41.3 & 76.0 & 31{,}242 & 151{,}814 \\
S-PR & 4.156 & 0.757 & 39.1 & 145.0 & 30{,}791 & 137{,}087 \\
\bottomrule
\end{tabular}
}
\end{table}


Most of the timed call is spent before the first iteration. The call indexes the nodes, assembles the sparse weighted adjacency matrix from the edge list in a per-edge Python loop, row-normalizes it and forms the restart vector, and only then runs power iteration; the edge lists themselves are built before the timer starts. A separate profiling run of the same call on the same edge lists attributes about 94\% of the wall time to matrix assembly ($0.52$ of $0.55$\,s for EndorseRank, $4.84$ of $5.17$\,s for AWP). The power iterations take about 1\%: about $6$\,ms for EndorseRank's $37$ iterations and $67$\,ms for AWP's $91$. Node indexing takes most of the rest, and forming AWP's restart vector about $11$\,ms. Assembly and each iteration are both linear in $|E|$ (Appendix~\ref{sec:pr-sparsity}), so the measured ratio of $9.2\times$ tracks the edge ratio of $9.3\times$. EndorseRank also converges in fewer iterations, $37$ against $91$. The allowance graph has no cycle apart from one self-loop, and $99.0\%$ of its edges end at spenders without out-edges, so every path stops after a step and the mass returns through the restart distribution. The transfer graph has a strongly connected core of $9{,}278$ addresses, and $28.5\%$ of its edges are reciprocated, so the walk circulates for longer before it settles. That difference is real, but it adds little to the wall time. Peak memory likewise grows with $|E|$, because the adjacency, weight and residual buffers are all proportional to the length of the edge list. Run-to-run noise is about two orders of magnitude smaller than the gap between the methods: the standard deviations are $0.014$\,s and $0.062$\,s.

The two hybrid rows show what the speed result covers and where it stops. C-PR walks on the union of both edge sets ($151{,}814$ edges, $31{,}242$ nodes) and needs $5.67$\,s, longer than AWP. It converges in fewer iterations ($76$ against $91$), but it has more edges to assemble, and assembly dominates the call. S-PR runs on the transfer layer and takes $5.60$\,s. Both its time and its iteration count include the solve of the allowance layer that supplies its restart vector: the $145$ iterations are $37$ for that solve plus $108$ for the seeded transfer walk. Claim~C1 is thus a statement about EndorseRank alone: any algorithm that queries the transfer layer, on its own or combined with allowance information, pays the cost of that query.

Two caveats attach to this result. The time covers only the solver call, matrix assembly included, on edge lists built in advance. A full end-to-end refresh would add parsing, decoding and edge-list construction, steps that all algorithms share wholly or in part, and C1 leaves them out for that reason. A faster method is also not necessarily a better one. The runtime result has to be read against the alignment results of Sections~\ref{sec:three-method-alignment}--\ref{sec:proxy-family-detail}, and against the holdout results that follow.

\dissertationSubsubheading[sub:runtime-scaling]{Node-count sensitivity on the expanded pool}

Deterministic SHA256 subsampling of the expanded pool yields three increasing node counts. Table~\ref{tab:benchmark-scaling} reports the runtime at each of them, together with the nodes and edges of every algorithm's subgraph at that stage. The in-cohort stages ($n\le 5{,}521$) are in Appendix~\ref{ch:appendix-eval} (Table~\ref{tab:benchmark-scaling-incohort}).

% Real BigQuery data -- regenerate with:
%   2-Dissertation-Draft-Works/wallet-reputation-experiments/scripts/run_dissertation_eval.py --real --export-latex
% Generated by export_latex_results.py -- do not edit by hand unless re-export fails
\begin{table}[htbp]
\centering
\caption[Runtime versus evaluation-set size on the expanded wallet pool.]{Runtime versus evaluation-set size on the expanded wallet pool ($N=99{,}080$; SHA256 subsampling seed \texttt{benchmark-tier2-v1}). Each stage keeps the edges that touch at least one sampled wallet, so $|V|$ and $|E|$ are reported per stage; the pool adds addresses but no additional extracted edges beyond the matched-cohort subgraph, hence the full-pool row reproduces the matched-cohort graph. Speedup = AWP/ER runtime ratio within the same row. Runtime = mean wall-clock seconds over 5 timed PageRank solves after one untimed warm-up; SD over the same 5 runs.}
\label{tab:benchmark-scaling}
\fitwidth{%
\begin{tabular}{rrrrrrrr}
\toprule
$n$ & ER (s) & AWP (s) & Speedup & ER $|V|$ & ER $|E|$ & AWP $|V|$ & AWP $|E|$ \\
\midrule
10{,}000 & 0.070 & 0.854 & 12.2$\times$ & 1{,}296 & 2{,}055 & 8{,}489 & 26{,}732 \\
50{,}000 & 0.334 & 2.706 & 8.1$\times$ & 4{,}440 & 9{,}190 & 19{,}990 & 92{,}887 \\
99{,}080 & 0.401 & 4.105 & 10.2$\times$ & 6{,}442 & 14{,}727 & 30{,}791 & 137{,}087 \\
\bottomrule
\end{tabular}
}
\end{table}


Two features of Table~\ref{tab:benchmark-scaling} need to be read together with Section~\ref{sub:expanded-pool}. The first is that the full-pool row is identical to the matched-cohort row of Table~\ref{tab:benchmark-runtime}. This is by design and is not an error. No approval or transfer logs were parsed for the extra addresses, so the subgraph induced at $N=99{,}080$ is exactly the cohort subgraph; our benchmarking procedure times each distinct graph once per session and reuses that timing whenever the same graph is tabulated again. The same measurement appears in the $d=0.85$ row of Table~\ref{tab:robustness-damping}. The second feature concerns $|E|$. Within each scaling table it grows with $n$, but not across the two samplings: the pool stage with $n=10{,}000$ has fewer EndorseRank edges ($2{,}055$) than the in-cohort stage with $n=1{,}000$ ($2{,}584$, Table~\ref{tab:benchmark-scaling-incohort}). Each stage keeps the extracted edges that touch the wallets the hash function selects. Only the cohort wallets, about 6\% of the pool, had their logs extracted, so most addresses in a pool sample bring no edges of their own. What the table does support is a bounded claim. Across stages whose edge counts differ by a factor of five to seven, AWP takes $8.6$ to $14.4$ times as long as EndorseRank, and EndorseRank stays at or below $0.55$\,s. The largest ratio belongs to the smallest pool stage, where AWP needs $97$ iterations and EndorseRank $32$. How either algorithm scales as edges grow with the number of addresses is not shown, because the second-stage extraction that was to supply the extra edges did not finish. Figure~\ref{fig:runtime-scaling} plots the in-cohort and pool stages on one log axis.

\begin{figure}[htbp]
\centering
\includegraphics[width=0.9\textwidth]{\figpath{runtime-scaling.pdf}}
\caption[{PageRank wall-clock runtime vs.\ number of sampled wallets (log $x$-axis). Solid lines: in-cohort stages ($n\le 5{,}521$, Table~\ref{tab:benchmark-scaling-incohort}). Dashed lines: expanded-pool stages ($N=99{,}080$, Table~\ref{tab:benchmark-scaling}). The pool stages add addresses but do not add any edges from extraction, so the rightmost pool point is at the full matched cohort.}]{Wall-clock time of one solver call against the number of sampled wallets (log $x$-axis). Solid lines: in-cohort stages ($n\le 5{,}521$, Table~\ref{tab:benchmark-scaling-incohort}). Dashed lines: expanded-pool stages ($N=99{,}080$, Table~\ref{tab:benchmark-scaling}). Pool stages add addresses without adding extracted edges, which puts the rightmost pool point at the full matched cohort. A call assembles the sparse transition matrix from the edge list and then iterates; assembly takes almost all of the time, so runtime follows the edge count $|E|$ of each stage and not $n$.}
\label{fig:runtime-scaling}
\end{figure}

\dissertationSubheading[sec:three-method-alignment]{Same-window alignment: EndorseRank vs.\ AWP (Claims~C2--C3)}

Claim~C2 asks whether each method recovers its own construct on same-window local statistics. Claim~C3 asks whether the winner is construct-specific, that is, whether the better of the two methods changes from one family to another. Table~\ref{tab:alignment-family-ci} reports the family-mean Kendall $\tau$ for the transfer, allowance and Sybil-stability families, each with a 95\% paired-bootstrap confidence interval, together with the inter-method agreement between the two score vectors. The same family means are drawn as a heatmap in Figure~\ref{fig:alignment-heatmap}. None of the comparisons below was fixed in advance (Chapter~\ref{ch:methodology}).

% Real BigQuery data -- regenerate with:
%   2-Dissertation-Draft-Works/wallet-reputation-experiments/scripts/run_dissertation_eval.py --real --export-latex
% Generated by export_latex_results.py -- do not edit by hand unless re-export fails
\begin{table}[htbp]
\centering
\caption[Family-mean Kendall $\tau$ with bootstrap confidence intervals on the matched cohort.]{Family-mean Kendall $\tau$ with bootstrap confidence intervals on the matched cohort ($n=5{,}521$). CI = 95\% percentile bootstrap over wallets (400 paired resamples, seed 42).}
\label{tab:alignment-family-ci}
\fitwidth{%
\begin{tabular}{lcccc}
\toprule
Family & ER mean $\tau$ & ER 95\% CI & AWP mean $\tau$ & AWP 95\% CI \\
\midrule
Transfer & 0.333 & [0.314, 0.353] & 0.534 & [0.522, 0.547] \\
Allowance & 0.355 & [0.322, 0.384] & -0.030 & [-0.045, -0.017] \\
Sybil & 0.233 & [0.214, 0.253] & 0.286 & [0.274, 0.300] \\
\midrule
\multicolumn{5}{l}{EndorseRank vs.\ AWP score agreement: $\tau=0.316$, CI [0.298, 0.336]} \\
\bottomrule
\end{tabular}
}
\end{table}


\begin{figure}[htbp]
\centering
\includegraphics[width=0.8\textwidth]{\figpath{alignment-heatmap.pdf}}
\caption[{Family-mean Kendall $\tau$ on the matched cohort ($n=5{,}521$) for EndorseRank and AWP; these are the values of Table~\ref{tab:alignment-family-ci}.}]{Family-mean Kendall $\tau$ for EndorseRank and AWP on the matched cohort ($n=5{,}521$), plotted from the values in Table~\ref{tab:alignment-family-ci}. The color scale is centered at zero, so red cells are positive and blue cells negative.}
\label{fig:alignment-heatmap}
\end{figure}

The two patterns cross. AWP tracks the transfer family most closely ($0.612$) and is slightly negative on the allowance family, while EndorseRank tracks the allowance family most closely ($0.375$) and the transfer family moderately ($0.330$). Table~\ref{tab:tau-diff} tests the five paired differences of Section~\ref{sec:statistical-tests}. The hybrid contrasts are reported in Table~\ref{tab:tau-diff-hybrid} (Section~\ref{sec:coupled-results}).

% Real BigQuery data -- regenerate with:
%   A-Skill-Programs/margin_rank/scripts/run_dissertation_eval.py --real --export-latex
% Generated by export_latex_results.py -- do not edit by hand unless re-export fails
\begin{table}[htbp]
\centering
\caption[Paired bootstrap tests of differences between Kendall $\tau$ coefficients.]{Paired bootstrap tests of differences between Kendall $\tau$ coefficients (matched cohort, $n=5{,}521$). Family entries use the family-mean $\tau$; the same wallet resamples are used for $a$ and $b$, so the interval is on the paired difference. An interval excluding 0 indicates a difference not attributable to sampling variation at the 5\% level. CI = 95\% percentile bootstrap over wallets (400 paired resamples, seed 42).}
\label{tab:tau-diff}
\fitwidth{%
\begin{tabular}{lccccc}
\toprule
Contrast ($a$ minus $b$) & $\tau_a$ & $\tau_b$ & $\Delta\tau$ & 95\% CI of $\Delta\tau$ & $P(\Delta\tau>0)$ \\
\midrule
EndorseRank allowance minus EndorseRank transfer & 0.355 & 0.333 & +0.022 & [-0.026, 0.067] & 0.810 \\
AWP transfer minus EndorseRank transfer & 0.534 & 0.333 & +0.201 & [0.182, 0.218] & 1.000 \\
EndorseRank allowance minus AWP allowance & 0.355 & -0.030 & +0.385 & [0.350, 0.421] & 1.000 \\
AWP sybil-stability minus EndorseRank sybil-stability & 0.286 & 0.233 & +0.052 & [0.038, 0.065] & 1.000 \\
EndorseRank in-approve degree minus EndorseRank in-degree & 0.356 & 0.344 & +0.012 & [-0.041, 0.060] & 0.675 \\
\bottomrule
\end{tabular}
}
\end{table}


In three of the five contrasts the interval excludes zero. AWP's transfer alignment is higher than EndorseRank's ($+0.282$), EndorseRank's allowance alignment is higher than AWP's ($+0.398$), and AWP leads EndorseRank on the Sybil-stability family ($+0.047$). The remaining two are null: EndorseRank's allowance alignment is not shown to be higher than its own transfer alignment (contrast~i, $+0.045$, interval $[-0.001,0.084]$), and its in-approve-degree alignment is not shown to be higher than its in-degree alignment (contrast~v, $+0.035$, $[-0.015,0.079]$). These two nulls cannot be read as evidence either way. EndorseRank ties most of the cohort (Section~\ref{sec:rank-divergence}), and the ties cap $\tau_b$ at different heights for different proxies: about $0.38$ for the allowance proxies, which EndorseRank reaches, and about $0.57$ for a proxy without ties. Coefficients with different ceilings are not on one scale. When the wallets outside the allowance graph are kept as isolated nodes instead of being scored zero, the ceilings move and both contrasts come out large and positive (Table~\ref{tab:tie-sensitivity}). Since EndorseRank smooths spender in-approve degree over the graph, its agreement with the allowance proxies is an intended-construct check and cannot count as an external test. On this cohort the check rests on the 131 wallets that receive approvals, all of which EndorseRank lifts above its common value. AWP spreads those 131 over most of its ranking, so transfer centrality says little about which of the cohort's wallets are approved as spenders. EndorseRank's transfer coefficient ($0.330$, interval $[0.311,0.350]$) is below AWP's and clear of zero, but most of the pairs behind it set a wallet at EndorseRank's common value against a zero-score wallet (Section~\ref{sec:rank-divergence}), and with those wallets kept as isolated nodes it turns negative.

Claim~C2 is supported insofar as each method reaches non-trivial alignment on its own construct, with an interval well clear of zero. The two between-method contrasts (ii and iii) support Claim~C3. Their intervals exclude zero in opposite directions, so neither method dominates both checks.

\dissertationSubheading[sec:proxy-family-detail]{Proxy family detail}

Each family mean condenses several per-proxy coefficients. This section reports those coefficients one proxy at a time, as Spearman $\rho$, Kendall $\tau$ and the 95\% interval of $\tau$. Before correlation, every proxy is oriented so that a higher value means higher expected reputation, and the transfer proxies are the ones AWP's source paper used to validate it\cend{3}.

\dissertationSubsubheading[sub:transfer-family]{Transfer family}

% Real BigQuery data -- regenerate with:
%   A-Skill-Programs/margin_rank/scripts/run_dissertation_eval.py --real --export-latex
% Generated by export_latex_results.py -- do not edit by hand unless re-export fails
\begin{table}[htbp]
\centering
\caption[Domain alignment with transfer proxies.]{Domain alignment with transfer proxies (inbound in-degree and in-value, the validation proxies used for the AWP baseline). Kendall $\tau$ and Spearman $\rho$ between reputation scores and proxy rankings. CI = 95\% percentile bootstrap over wallets (400 paired resamples, seed 42).}
\label{tab:alignment-transfer}
\fitwidth{%
\begin{tabular}{lccc}
\toprule
Proxy & Spearman $\rho$ & Kendall $\tau$ & 95\% CI of $\tau$ \\
\midrule
\multicolumn{4}{l}{\textit{EndorseRank}} \\
In-degree & 0.422 & 0.344 & [0.324, 0.366] \\
In-value (sum inbound) & 0.398 & 0.321 & [0.303, 0.342] \\
\addlinespace
\multicolumn{4}{l}{\textit{AWP}} \\
In-degree & 0.791 & 0.617 & [0.604, 0.631] \\
In-value (sum inbound) & 0.610 & 0.450 & [0.436, 0.467] \\

\bottomrule
\end{tabular}
}
\end{table}


In Table~\ref{tab:alignment-transfer}, AWP's agreement with in-degree ($\tau=0.728$) is the largest per-proxy Kendall $\tau$ on the matched cohort. Its agreement with in-value is much lower ($0.495$). That is what its weights lead one to expect: $V$ counts every transfer of ten base units or more about once, whatever its amount, so AWP follows the number of incoming transfers more closely than their summed value. The two are primary-construct checks for AWP, and AWP does not collapse to either; propagation, time weights and restarts at active senders shift a measurable share of pairwise orderings. EndorseRank's transfer coefficients are moderate ($0.341$ and $0.319$), and their intervals overlap neither zero nor AWP's intervals. They mainly reflect that wallets which made an approval in the window also receive more transfers than wallets which made none (Section~\ref{sec:rank-divergence}). A protocol that equates reputation with inbound flow can read that flow from the in-degree itself; AWP adds propagation and weighting, and out of window it ranks new senders below the in-degree (Section~\ref{sec:holdout-results}).

\dissertationSubsubheading[sub:allowance-family]{Allowance family}

% Real BigQuery data -- regenerate with:
%   1-Dissertation-Works/ERC20-Allowance-PageRank-Wallet-Reputation/dissertation/wallet-reputation-experiments/scripts/run_dissertation_eval.py --real --export-latex
% Generated by export_latex_results.py -- do not edit by hand unless re-export fails
\begin{table}[htbp]
\centering
\caption[Domain alignment with allowance proxies.]{Domain alignment with allowance proxies (spender-side in-approve degree and value). CI = 95\% percentile bootstrap over wallets (400 paired resamples, seed 42).}
\label{tab:alignment-allowance}
\fitwidth{%
\begin{tabular}{lccc}
\toprule
Proxy & Spearman $\rho$ & Kendall $\tau$ & 95\% CI of $\tau$ \\
\midrule
\multicolumn{4}{l}{\textit{EndorseRank}} \\
In-approve degree & 0.380 & 0.376 & [0.344, 0.403] \\
In-approve value (sum) & 0.380 & 0.373 & [0.341, 0.398] \\
\addlinespace
\multicolumn{4}{l}{\textit{AWP}} \\
In-approve degree & -0.029 & -0.024 & [-0.043, -0.002] \\
In-approve value (sum) & -0.029 & -0.023 & [-0.043, -0.002] \\

\bottomrule
\end{tabular}
}
\end{table}


In Table~\ref{tab:alignment-allowance}, EndorseRank's coefficients on in-approve degree and in-approve value are almost equal ($0.376$ and $0.373$). Both proxies are positive for the same 131 wallets, and EndorseRank lifts all of them above its common value. The two coefficients therefore mostly measure the same split, and they sit at the ceiling that the ties allow ($0.377$ and $0.378$). AWP's coefficients are slightly negative ($-0.024$ and $-0.023$), with intervals that exclude zero: among the cohort's wallets, transfer centrality does not pick out those that are approved as spenders. The comparison rests on these 131 wallets and is partly built in, since EndorseRank is computed from the same approvals. The same construct is therefore tested out of window on the spender holdout of Section~\ref{sec:holdout-results}, where every wallet has received approvals.

\dissertationSubsubheading[sub:sybil-family]{Sybil-adjusted stability family}

% Real BigQuery data -- regenerate with:
%   1-Dissertation-Works/ERC20-Allowance-PageRank-Wallet-Reputation/dissertation/wallet-reputation-experiments/scripts/run_dissertation_eval.py --real --export-latex
% Generated by export_latex_results.py -- do not edit by hand unless re-export fails
\begin{table}[htbp]
\centering
\caption[Domain alignment with Sybil-adjusted stability proxies from transfer activity.]{Domain alignment with Sybil-adjusted stability proxies from transfer activity. CI = 95\% percentile bootstrap over wallets (400 paired resamples, seed 42).}
\label{tab:alignment-sybil-stability}
\fitwidth{%
\begin{tabular}{lccc}
\toprule
Proxy & Spearman $\rho$ & Kendall $\tau$ & 95\% CI of $\tau$ \\
\midrule
\multicolumn{4}{l}{\textit{EndorseRank}} \\
Inbound counterparty ratio & 0.091 & 0.072 & [0.045, 0.095] \\
Transfer tenure (days) & 0.363 & 0.294 & [0.274, 0.316] \\
Active months & 0.377 & 0.327 & [0.307, 0.348] \\
\addlinespace
\multicolumn{4}{l}{\textit{AWP}} \\
Inbound counterparty ratio & -0.269 & -0.228 & [-0.250, -0.207] \\
Transfer tenure (days) & 0.685 & 0.507 & [0.493, 0.521] \\
Active months & 0.704 & 0.555 & [0.540, 0.569] \\

\bottomrule
\end{tabular}
}
\end{table}


Table~\ref{tab:alignment-sybil-stability} shows AWP ahead of EndorseRank on tenure ($0.507$ against $0.294$) and active months ($0.555$ against $0.327$), as one would expect from a transfer score that rewards activity spread over the window; the paired difference on the family mean excludes zero (Table~\ref{tab:tau-diff}, contrast~iv). On the inbound-counterparty ratio AWP is clearly negative ($-0.228$). Its edge weights count transfers, so a wallet that receives many transfers from a few senders ranks high, and that is the pattern the ratio penalizes. EndorseRank is weakly positive on the ratio ($0.072$). The family is reported as a robustness note on how much longitudinal transfer information each method carries, and not as a headline claim.

\dissertationSubheading[sec:rank-divergence]{Rank divergence (H1 / RQ1)}

Hypothesis~H1 and RQ1 concern whether EndorseRank and AWP put the same wallets in different orders. The inter-method Kendall $\tau_b$ between the two score vectors is $0.360$, with interval $[0.342,0.379]$ (last row of Table~\ref{tab:alignment-family-ci}), so the interval excludes both zero and one. The two rankings are related, yet far from interchangeable. Agreement runs above chance mainly because the wallets that EndorseRank scores zero also sit near the bottom of AWP (see below). Because $67.3\%$ of wallet pairs tie on EndorseRank, only $32.7\%$ of all pairs are ordered by both scores, and $18.6\%$ of those pairs are discordant.

The marginal score distributions are in Figure~\ref{fig:score-distributions}, and the EndorseRank distribution is dominated by one tie block. EndorseRank takes only $99$ distinct values on the cohort: $4{,}424$ wallets share the score $9.62\times10^{-5}$ that a wallet receiving no allowance gets from teleportation and redistributed dangling mass; $966$ score zero because no allowance edge touches them; only $131$ lie above the common value. Ties are rare under AWP, at $0.5\%$ of wallet pairs. A marginal distribution cannot show how two scores order the same wallets, so Figure~\ref{fig:rank-scatter} plots the two rankings against each other, with tied wallets at their mid-rank as in $\tau_b$. The EndorseRank tie blocks therefore appear as vertical columns. The $131$ wallets that EndorseRank separates are spread over most of the AWP ranking (median AWP rank $3{,}096$, interquartile range $1{,}520$--$4{,}850$), so its leaders are not AWP's leaders. The column of $4{,}424$ tied wallets spans almost the whole AWP range (interquartile range $1{,}149$--$3{,}571$). The $966$ zero-score wallets sit mostly near the bottom of AWP (median AWP rank $4{,}966$), and $370$ of them score zero under AWP as well.

The tie blocks come from two rules, and one of them is a choice of this pipeline. The 4,424 wallets at the common value are in the allowance graph as owners that receive no approval, so PageRank gives them teleportation and dangling mass and nothing else. The 966 zero-score wallets are not in the graph at all, and the pipeline scores a missing wallet zero (Section~\ref{sec:graph-construction}); kept as isolated nodes, they would have received the common value too. This rule decides most of EndorseRank's ordering of the cohort. Of the 4.99 million wallet pairs that EndorseRank does not tie, 4.27 million (85.7\%) set a wallet at the common value against a zero-score wallet. On the matched cohort, EndorseRank therefore mostly separates wallets that made an approval in the window from wallets that made none. The ties also cap $\tau_b$. Against the allowance proxies, which are positive for the 131 wallets above the common value, the largest $\tau_b$ that EndorseRank can reach with these ties is about $0.38$; against a proxy without ties the ceiling is about $0.57$.

Nothing in the graph favors the zero rule over the isolated-node alternative, so we reran the same-window alignment with the missing wallets kept as isolated nodes. This check was run after the results above were known. Under uniform teleportation an isolated node scores like any node without in-edges, so the 966 wallets join the block of 4,424 and every other score is multiplied by one common factor. AWP's 381 missing wallets were handled the same way; under AWP's restarts an isolated wallet that has sent nothing receives no restart mass, so it scores zero either way. Table~\ref{tab:tie-sensitivity} gives the result. No AWP coefficient moves. EndorseRank's allowance coefficient rises to $0.990$ (interval $[0.988,0.992]$), since EndorseRank and the allowance proxies now split the cohort at the same place. Its transfer coefficient falls to $-0.071$ ($[-0.086,-0.056]$), and its Sybil-stability and trading-success coefficients to $-0.036$ and $-0.017$. The agreement between the two scores falls from $0.360$ to $-0.025$ ($[-0.044,-0.003]$). Contrasts~i and~v, null under the zero rule, become $+1.061$ and $+1.090$, with intervals far above zero. EndorseRank's positive same-window coefficients on transfer, stability and trading, and its agreement with AWP, therefore measure which traders made an approval in the window.

% Post hoc checks -- regenerate with:
%   A-Skill-Programs/margin_rank/scripts/run_supplementary_checks.py
% Generated by export_latex_results.py -- do not edit by hand unless re-export fails
\begin{table}[htbp]
\centering
\caption[Same-window family-mean Kendall $\tau$ with wallets missing from a graph scored zero or kept as isolated nodes.]{Same-window family-mean Kendall $\tau$ with wallets missing from a graph scored zero or kept as isolated nodes (matched cohort, $n=5{,}521$). The main analysis scores a wallet that is missing from a method's graph zero; as an isolated node it gets the score of a node without in-edges. The rule concerns 966 wallets for EndorseRank and 381 for AWP. The last row is the agreement between the two scores. The check was run after the results were known. CI = 95\% percentile bootstrap over wallets (400 paired resamples, seed 42).}
\label{tab:tie-sensitivity}
\fitwidth{%
\begin{tabular}{lcccc}
\toprule
 & \multicolumn{2}{c}{Missing wallets scored zero} & \multicolumn{2}{c}{Missing wallets as isolated nodes} \\
\cmidrule(lr){2-3} \cmidrule(lr){4-5}
Family & EndorseRank & AWP & EndorseRank & AWP \\
\midrule
Transfer & 0.333 [0.314, 0.353] & 0.534 [0.522, 0.547] & -0.071 [-0.086, -0.057] & 0.534 [0.522, 0.547] \\
Allowance & 0.355 [0.322, 0.384] & -0.030 [-0.045, -0.017] & 0.956 [0.928, 0.976] & -0.030 [-0.045, -0.017] \\
Sybil stability & 0.233 [0.214, 0.253] & 0.286 [0.274, 0.300] & -0.033 [-0.046, -0.020] & 0.285 [0.273, 0.300] \\
\addlinespace
Liquidation & 0.093 [0.078, 0.110] & 0.113 [0.098, 0.128] & 0.082 [0.068, 0.095] & 0.113 [0.098, 0.128] \\
Inverse risk & -0.056 [-0.072, -0.040] & -0.094 [-0.108, -0.081] & 0.003 [-0.019, 0.019] & -0.094 [-0.108, -0.081] \\
GMX trading success & 0.096 [0.080, 0.113] & 0.179 [0.166, 0.193] & -0.019 [-0.034, -0.004] & 0.179 [0.166, 0.193] \\
\midrule
EndorseRank vs.\ AWP & \multicolumn{2}{c}{0.316 [0.298, 0.336]} & \multicolumn{2}{c}{-0.025 [-0.038, -0.012]} \\
\bottomrule
\end{tabular}
}
\end{table}


\begin{figure}[htbp]
\centering
\includegraphics[width=0.9\textwidth]{\figpath{score-distributions.pdf}}
\caption[{Marginal distributions of EndorseRank and AWP scores on the matched cohort ($n=5{,}521$), on a $\log_{10}$ scale for wallets with positive score (EndorseRank has $966$ zero-score wallets and AWP $381$, omitted from the histograms).}]{Marginal EndorseRank and AWP score distributions on the matched cohort ($n=5{,}521$), drawn on a $\log_{10}$ scale for wallets with a positive score; counts are on a log scale. Zero-score wallets, which no edge of the respective graph touches, are left out: $966$ for EndorseRank and $381$ for AWP. For EndorseRank, $4{,}424$ wallets share the score $9.62\times10^{-5}$ of a wallet that receives no allowance, and only $131$ lie above it.}
\label{fig:score-distributions}
\end{figure}

\begin{figure}[htbp]
\centering
\includegraphics[width=0.8\textwidth]{\figpath{rank-scatter.pdf}}
\caption[{Rank--rank scatter plot of EndorseRank vs.\ AWP for the matched cohort (random sample of $2{,}000$ wallets). Points on the dashed identity line represent identical orderings.}]{Rank--rank scatter of EndorseRank against AWP on the matched cohort (random sample of $2{,}000$ wallets). Tied scores share their average rank, which is the tie treatment of Kendall's $\tau_b$. The two vertical columns are EndorseRank tie blocks: the $4{,}424$ wallets with the common score $9.62\times10^{-5}$ and the $966$ wallets with score zero. A point on the dashed identity line has the same rank in both orderings.}
\label{fig:rank-scatter}
\end{figure}

Construct-specific alignment has a structural counterpart in rank divergence. Had the two methods ordered wallets almost identically, the family-level winners in Table~\ref{tab:alignment-family-ci} might be noise around one shared ranking. The divergence actually observed, read together with the opposite-signed between-method contrasts of Table~\ref{tab:tau-diff}, supports treating the two as distinct instruments, related mainly through the wallets that EndorseRank scores zero. H1 is supported. On this cohort, much of the difference comes from EndorseRank's ties, so the claim that the two scores measure different things rests mainly on the holdout (Section~\ref{sec:holdout-results}). In operational terms, a protocol now relying on transfer-based scores cannot assume that its wallet ordering will survive a switch to allowance-based scores, or a blend with them. The largest disagreements are expected among transfer hubs rarely endorsed as spenders and among endorsed spenders with small transfer footprints.

\dissertationSubheading[sec:robustness-results]{Robustness and sensitivity checks}

The results above would be of limited use if they hinged on one damping value, one token mix or one sample definition. Three checks address this. The first two are robustness checks in the strict sense, holding the sample fixed while a parameter or the token universe is perturbed. The third is a sensitivity analysis: it varies the sample definition itself, and there we expect the coefficients to move. A fourth check asks how much EndorseRank depends on its restart rule.

\dissertationSubsubheading[sub:damping]{Damping}

% Real BigQuery data -- regenerate with:
%   2-Dissertation-Draft-Works/wallet-reputation-experiments/scripts/run_dissertation_eval.py --real --export-latex
% Generated by export_latex_results.py -- do not edit by hand unless re-export fails
\begin{table}[htbp]
\centering
\caption[Robustness: mean Kendall $\tau$ under PageRank damping $d \in \{0.75, 0.85, 0.95\}$ on the matched alignment cohort.]{Robustness: mean Kendall $\tau$ under PageRank damping $d \in \{0.75, 0.85, 0.95\}$ on the matched alignment cohort ($n=5{,}521$). The $d=0.85$ runtimes are the same measurements as Table~\ref{tab:benchmark-runtime} (one timing session; each graph--damping pair is timed once).}
\label{tab:robustness-damping}
\fitwidth{%
\begin{tabular}{rcccccc}
\toprule
$d$ & ER all. & AWP all. & ER tr. & AWP tr. & ER (s) & AWP (s) \\
\midrule
0.75 & 0.355 & -0.031 & 0.333 & 0.531 & 0.594 & 4.420 \\
0.85 & 0.355 & -0.030 & 0.333 & 0.534 & 0.401 & 4.105 \\
0.95 & 0.355 & -0.030 & 0.333 & 0.536 & 0.427 & 5.420 \\
\bottomrule
\end{tabular}
}
\end{table}


For EndorseRank the family means do not move at the third decimal place across $d\in\{0.75,0.85,0.95\}$. For AWP the transfer coefficient falls from $0.626$ to $0.585$ as $d$ rises, and the allowance coefficient stays slightly negative. The crossing of the patterns is no artifact of the conventional $d=0.85$. EndorseRank's stability is largely structural. Only $131$ cohort wallets lie above EndorseRank's common value (Section~\ref{sec:rank-divergence}); a change in $d$ moves the common value but leaves the $4{,}424$ wallets that share it tied, so the rank correlations can hardly move. Runtime barely moves either. EndorseRank takes $0.541$, $0.547$ and $0.541$\,s at $d=0.75$, $0.85$ and $0.95$, and AWP $5.011$, $5.029$ and $5.187$\,s, so the ratio stays between $9.2\times$ and $9.6\times$. Each graph--damping pair was timed once, in one session. Matrix assembly takes most of each call and does not depend on $d$, which is why the extra iterations at $d=0.95$ add so little. Figure~\ref{fig:damping-sensitivity} plots the family means.

\begin{figure}[htbp]
\centering
\includegraphics[width=0.85\textwidth]{\figpath{damping-sensitivity.pdf}}
\caption[{Sensitivity of family-mean Kendall $\tau$ for the allowance and transfer families over $d\in\{0.75,0.85,0.95\}$ (Table~\ref{tab:robustness-damping} values).}]{Damping sensitivity of family-mean Kendall $\tau$ for the allowance and transfer families across $d\in\{0.75,0.85,0.95\}$ (values from Table~\ref{tab:robustness-damping}).}
\label{fig:damping-sensitivity}
\end{figure}

\dissertationSubsubheading[sub:token-subgraph]{Top-token subgraph}

% Real BigQuery data -- regenerate with:
%   A-Skill-Programs/margin_rank/scripts/run_dissertation_eval.py --real --export-latex
% Generated by export_latex_results.py -- do not edit by hand unless re-export fails
\begin{table}[htbp]
\centering
\caption{Robustness: mean Kendall $\tau$ on the matched cohort when both graphs keep only 20 ERC-20 tokens, chosen by summed raw amount or by number of rows (transfer events and latest allowances), next to the value on all tokens from Table~\ref{tab:alignment-family-ci}.}
\label{tab:robustness-tokens}
\fitwidth{%
\begin{tabular}{lccccccccc}
\toprule
Method & \multicolumn{3}{c}{Transfer} & \multicolumn{3}{c}{Allowance} & \multicolumn{3}{c}{Sybil} \\
\cmidrule(lr){2-4} \cmidrule(lr){5-7} \cmidrule(lr){8-10}
 & by amount & by count & all & by amount & by count & all & by amount & by count & all \\
\midrule
EndorseRank & 0.331 & 0.332 & 0.333 & 0.355 & 0.353 & 0.355 & 0.234 & 0.233 & 0.233 \\
AWP & 0.532 & 0.530 & 0.534 & -0.031 & -0.030 & -0.030 & 0.283 & 0.280 & 0.286 \\
\bottomrule
\end{tabular}
}
\end{table}


When both graphs are restricted to twenty tokens (Table~\ref{tab:robustness-tokens}), the largest change in a family mean is $0.014$, for AWP on the transfer family when the tokens are the twenty with the largest summed raw amounts; with the twenty tokens that have the most rows it is $0.009$. EndorseRank moves by $0.005$ at most. The two lists share eleven tokens, so the pattern does not depend on the many tokens outside either set. Neither list ranks tokens by economic value: summed raw amounts favor tokens with many unlimited approvals, and row counts favor tokens that change hands often.

\dissertationSubsubheading[sub:sample-size]{Sample-definition sensitivity}

% Real BigQuery data -- regenerate with:
%   2-Dissertation-Draft-Works/wallet-reputation-experiments/scripts/run_dissertation_eval.py --real --export-latex
% Generated by export_latex_results.py -- do not edit by hand unless re-export fails
\begin{table}[htbp]
\centering
\caption[Sensitivity of alignment to the evaluation set: mean Kendall $\tau$ and runtime at staged subsample sizes drawn from the expanded wallet pool.]{Sensitivity of alignment to the evaluation set: mean Kendall $\tau$ and runtime at staged subsample sizes drawn from the expanded wallet pool ($N=99{,}080$; seed \texttt{benchmark-tier2-v1}). Rows are evaluated on different wallet sets from the matched cohort, so $\tau$ values are not comparable with Table~\ref{tab:alignment-family-ci}; they show how much the statistic moves with the evaluation population. Runtimes are the same measurements as Table~\ref{tab:benchmark-scaling}.}
\label{tab:robustness-sample-size}
\fitwidth{%
\begin{tabular}{rccccrrc}
\toprule
$n$ & ER all. & AWP all. & ER tr. & AWP tr. & ER (s) & AWP (s) & Speedup \\
\midrule
10{,}000 & 0.343 & 0.242 & 0.821 & 0.904 & 0.070 & 0.854 & 12.20 \\
50{,}000 & 0.330 & 0.232 & 0.828 & 0.915 & 0.334 & 2.706 & 8.10 \\
99{,}080 & 0.327 & 0.230 & 0.828 & 0.916 & 0.401 & 4.105 & 10.24 \\
\bottomrule
\end{tabular}
}
\end{table}


In Table~\ref{tab:robustness-sample-size} the transfer and allowance family means are recomputed on the staged pool subsamples from the runtime scaling. The transfer-family means come out much larger than on the matched cohort ($0.82$--$0.91$), and AWP's allowance coefficient turns positive ($0.23$--$0.24$), while EndorseRank's allowance coefficient ($0.332$--$0.343$) ends below its cohort value of $0.375$. None of these shifts is evidence about the methods. The pool subsamples let in wallets with no edges in one graph or in both; after alignment such a wallet has a zero score and a zero proxy, and tied zeros at the bottom of both rankings push Kendall $\tau$ up mechanically (Section~\ref{sub:sample-size-protocol}). We therefore report the table as a sensitivity analysis of the sample definition. The question it answers is narrower, namely whether the relative ordering of the two methods within each family survives a change of population. At every stage it does, with AWP ahead on transfer and EndorseRank ahead on allowance. The absolute values, though, must not be compared with those of Table~\ref{tab:alignment-family-ci}. At every pool stage the EndorseRank graph has far fewer nodes than the sample has wallets (Table~\ref{tab:benchmark-scaling}), so most pool wallets score zero on it, against $966$ of the $5{,}521$ cohort wallets.

For AWP, the damping and top-token checks show that Claims~C2 and~C3 do not hinge on the parameter or the token mix; for EndorseRank they could hardly have moved, because its ties hold whatever $d$ or token set is used. The runtime ratio behind Claim~C1 stays between $9.2\times$ and $9.6\times$ across the damping sweep. When the population changes, as in the sample-definition check, the comparative structure survives but the absolute values do not.

\dissertationSubsubheading[sub:endorserank-restarts]{Restarts of EndorseRank}

EndorseRank restarts the walk uniformly, while AWP restarts it at addresses in proportion to their recent sending activity (Section~\ref{sub:endorserank-graph}). Carried over to allowances, AWP's rule restarts the walk at owners in proportion to the number of spenders they approved recently. Table~\ref{tab:endorserank-restarts} scores both versions of EndorseRank, with AWP beside them. The restart rule was chosen after this comparison had been made.

% Real BigQuery data -- regenerate with:
%   2-Dissertation-Draft-Works/wallet-reputation-experiments/scripts/run_dissertation_eval.py --real --export-latex
% Generated by export_latex_results.py -- do not edit by hand unless re-export fails
\begin{table}[htbp]
\centering
\caption[EndorseRank with uniform restarts and with restarts weighted by approving activity, AWP's rule applied to the allowance graph, next to AWP: family-mean Kendall $\tau$ on the matched cohort ($n=5{,}521$) and $\tau$ with the two spring labels on the spender cohort.]{EndorseRank with uniform restarts and with restarts weighted by approving activity, AWP's rule applied to the allowance graph, next to AWP: family-mean Kendall $\tau$ on the matched cohort ($n=5{,}521$) and $\tau$ with the two spring labels on the spender cohort ($n=1{,}335$). Both EndorseRank rows weight each latest allowance by $\sigma(\Delta t)V(z)$. CI = 95\% percentile bootstrap over wallets (400 paired resamples, seed 42).}
\label{tab:endorserank-restarts}
\fitwidth{%
\begin{tabular}{llccccc}
\toprule
 & & \multicolumn{3}{c}{Same window, matched cohort} & \multicolumn{2}{c}{Spring holdout, spenders} \\
\cmidrule(lr){3-5} \cmidrule(lr){6-7}
Method & Restarts & Transfer & Allowance & Sybil & New approval pairs & New senders \\
\midrule
EndorseRank & uniform & \begin{tabular}[c]{@{}c@{}}0.330\\[-1pt]{\footnotesize [0.311, 0.350]}\end{tabular} & \begin{tabular}[c]{@{}c@{}}0.375\\[-1pt]{\footnotesize [0.342, 0.401]}\end{tabular} & \begin{tabular}[c]{@{}c@{}}0.231\\[-1pt]{\footnotesize [0.212, 0.250]}\end{tabular} & \begin{tabular}[c]{@{}c@{}}0.394\\[-1pt]{\footnotesize [0.360, 0.431]}\end{tabular} & \begin{tabular}[c]{@{}c@{}}0.283\\[-1pt]{\footnotesize [0.247, 0.320]}\end{tabular} \\
EndorseRank & by approving activity & \begin{tabular}[c]{@{}c@{}}0.445\\[-1pt]{\footnotesize [0.433, 0.457]}\end{tabular} & \begin{tabular}[c]{@{}c@{}}0.044\\[-1pt]{\footnotesize [0.029, 0.062]}\end{tabular} & \begin{tabular}[c]{@{}c@{}}0.245\\[-1pt]{\footnotesize [0.233, 0.259]}\end{tabular} & \begin{tabular}[c]{@{}c@{}}0.372\\[-1pt]{\footnotesize [0.327, 0.417]}\end{tabular} & \begin{tabular}[c]{@{}c@{}}0.301\\[-1pt]{\footnotesize [0.254, 0.348]}\end{tabular} \\
AWP & by sending activity & \begin{tabular}[c]{@{}c@{}}0.612\\[-1pt]{\footnotesize [0.602, 0.622]}\end{tabular} & \begin{tabular}[c]{@{}c@{}}-0.024\\[-1pt]{\footnotesize [-0.043, -0.002]}\end{tabular} & \begin{tabular}[c]{@{}c@{}}0.278\\[-1pt]{\footnotesize [0.264, 0.292]}\end{tabular} & \begin{tabular}[c]{@{}c@{}}0.123\\[-1pt]{\footnotesize [0.068, 0.177]}\end{tabular} & \begin{tabular}[c]{@{}c@{}}0.344\\[-1pt]{\footnotesize [0.307, 0.383]}\end{tabular} \\
\bottomrule
\end{tabular}
}
\end{table}


With AWP's restarts, EndorseRank's same-window allowance coefficient falls from $0.375$ to $0.044$, and its transfer coefficient rises from $0.330$ to $0.445$. The restart mass then goes to the owners that approve most. The cohort's traders are such owners, so the score orders them by how much they approve, which is closer to their activity than to the endorsements they receive. Out of window the two versions are close. On the spring spenders they reach $0.394$ and $0.372$ on new approval pairs and $0.283$ and $0.301$ on new senders, with overlapping intervals, and in June--August $0.333$ and $0.325$, and $0.255$ and $0.284$ (Table~\ref{tab:fresh-posthoc}). The uniform rule keeps the allowance construct in the same window and gives up nothing measurable out of window.

\dissertationSubheading[sec:coupled-results]{Coupled operator: same-window results (RQ5)}

Two hybrids were defined in Section~\ref{sub:coupled-graph}. The coupled operator C-PR mixes the raw-amount allowance and transfer layers in each wallet's transition row with weight $\lambda$, and the endorsement-seeded walk S-PR keeps the transfer layer but restarts from the allowance walk. Table~\ref{tab:alignment-hybrid} sets both beside EndorseRank, AWP and the two layers walked alone, C-PR ($\lambda=1$) and C-PR ($\lambda=0$), on the three same-window families of Table~\ref{tab:alignment-family-ci}. Table~\ref{tab:tau-diff-hybrid} gives the paired contrasts: first the eight of Section~\ref{sec:statistical-tests} against whichever layer leads each family, then four against EndorseRank and AWP. RQ5 is decided by the out-of-window tests in Section~\ref{sec:holdout-results}. Here we only locate the hybrids on the checks that the single-layer methods were built to pass.

% Real BigQuery data -- regenerate with:
%   2-Dissertation-Draft-Works/wallet-reputation-experiments/scripts/run_dissertation_eval.py --real --export-latex
% Generated by export_latex_results.py -- do not edit by hand unless re-export fails
\begin{table}[htbp]
\centering
\caption[Same-window family-mean Kendall $\tau$ with bootstrap intervals for the two single-layer methods, the coupled operator at three values of $\lambda$ and the endorsement-seeded walk.]{Same-window family-mean Kendall $\tau$ with bootstrap intervals for the two single-layer methods, the coupled operator at three values of $\lambda$ and the endorsement-seeded walk (matched cohort, $n=5{,}521$). $\lambda=1$ would reproduce EndorseRank and $\lambda=0$ AWP on a shared node set. CI = 95\% percentile bootstrap over wallets (400 paired resamples, seed 42).}
\label{tab:alignment-hybrid}
\fitwidth{%
\begin{tabular}{lccc}
\toprule
Method & Transfer & Allowance & Sybil \\
\midrule
EndorseRank & 0.333 [0.314, 0.353] & 0.355 [0.322, 0.384] & 0.233 [0.214, 0.253] \\
AWP & 0.534 [0.522, 0.547] & -0.030 [-0.045, -0.017] & 0.286 [0.274, 0.300] \\
\midrule
C-PR ($\lambda=0.25$) & 0.531 [0.520, 0.545] & -0.009 [-0.020, 0.004] & 0.288 [0.276, 0.302] \\
C-PR ($\lambda=0.5$) & 0.527 [0.515, 0.540] & -0.001 [-0.013, 0.013] & 0.290 [0.279, 0.304] \\
C-PR ($\lambda=0.75$) & 0.519 [0.507, 0.533] & 0.006 [-0.007, 0.020] & 0.292 [0.281, 0.306] \\
S-PR & 0.494 [0.481, 0.508] & 0.116 [0.097, 0.134] & 0.220 [0.207, 0.235] \\
\bottomrule
\end{tabular}
}
\end{table}


% Real BigQuery data -- regenerate with:
%   A-Skill-Programs/margin_rank/scripts/run_dissertation_eval.py --real --export-latex
% Generated by export_latex_results.py -- do not edit by hand unless re-export fails
\begin{table}[htbp]
\centering
\caption[Same-window paired contrasts for the coupled operator (C-PR, $\lambda=0.5$) and the endorsement-seeded walk (S-PR) against the single-layer method that leads each family.]{Same-window paired contrasts for the coupled operator (C-PR, $\lambda=0.5$) and the endorsement-seeded walk (S-PR) against the single-layer method that leads each family (matched cohort, $n=5{,}521$). Family-mean $\tau$; the same wallet resamples are used for $a$ and $b$. The secondary criterion of Section~\ref{sec:holdout-protocol} is met on a family when the interval includes 0 or lies above it (transfer, allowance) or lies above 0 (Sybil stability). CI = 95\% percentile bootstrap over wallets (400 paired resamples, seed 42).}
\label{tab:tau-diff-hybrid}
\fitwidth{%
\begin{tabular}{lccccc}
\toprule
Contrast ($a$ minus $b$) & $\tau_a$ & $\tau_b$ & $\Delta\tau$ & 95\% CI of $\Delta\tau$ & $P(\Delta\tau>0)$ \\
\midrule
C-PR transfer minus AWP transfer & 0.527 & 0.534 & -0.007 & [-0.009, -0.005] & 0.000 \\
C-PR allowance minus EndorseRank allowance & -0.001 & 0.355 & -0.356 & [-0.388, -0.323] & 0.000 \\
C-PR sybil-stability minus AWP sybil-stability & 0.290 & 0.286 & +0.004 & [0.003, 0.006] & 1.000 \\
C-PR sybil-stability minus EndorseRank sybil-stability & 0.290 & 0.233 & +0.057 & [0.043, 0.069] & 1.000 \\
S-PR transfer minus AWP transfer & 0.494 & 0.534 & -0.040 & [-0.049, -0.030] & 0.000 \\
S-PR allowance minus EndorseRank allowance & 0.116 & 0.355 & -0.239 & [-0.263, -0.215] & 0.000 \\
S-PR sybil-stability minus AWP sybil-stability & 0.220 & 0.286 & -0.065 & [-0.072, -0.058] & 0.000 \\
S-PR sybil-stability minus EndorseRank sybil-stability & 0.220 & 0.233 & -0.013 & [-0.025, -0.003] & 0.007 \\
\bottomrule
\end{tabular}
}
\end{table}


C-PR behaves like its transfer layer with a small allowance-layer correction. Its transfer-family coefficient falls from $0.531$ at $\lambda=0.25$ to $0.519$ at $\lambda=0.75$, always below C-PR ($\lambda=0$) at $0.534$ and within $0.015$ of it. At the primary setting the paired difference to that walk is $-0.007$ (interval $[-0.009,-0.005]$), small but clear of zero. AWP ranks the transfer family better than either: C-PR lies $0.085$ below it ($[-0.095,-0.075]$). On the allowance family C-PR cannot be told apart from zero at any $\lambda$ and lies far below its allowance layer ($\Delta\tau=-0.356$, interval $[-0.388,-0.323]$) and below EndorseRank ($-0.375$, $[-0.404,-0.344]$). The cause is the sparsity of the allowance layer. The transfer layer has about nine times the allowance layer's edges, and $30{,}791$ nodes against $6{,}442$ (Table~\ref{tab:benchmark-runtime}), so most nodes of the union graph have no allowance out-edge. At those nodes the mixing rule of Equation~\eqref{eq:coupled-transition} falls back on the transfer row whatever $\lambda$ is; $\lambda$ acts only at the minority of nodes with out-edges in both layers. Once the walk leaves that minority, it travels on transfer edges until it teleports. Raising $\lambda$ to $0.75$ shifts the allowance coefficient by less than one hundredth. The mixture recovers the allowance construct only at the end point $\lambda=1$; at $\lambda=0.75$ its allowance coefficient is still $0.006$.

Only on the Sybil-adjusted stability family does C-PR exceed both layers walked alone, as the secondary criterion of Section~\ref{sec:holdout-protocol} asked. It scores $0.290$, against $0.286$ for C-PR ($\lambda=0$) and $0.233$ for C-PR ($\lambda=1$), and both paired intervals lie above zero ($[0.003,0.006]$ and $[0.043,0.069]$). It also exceeds AWP ($0.278$) and EndorseRank ($0.231$) on this family. The gain over the transfer walk is four thousandths, which makes it a difference in the statistical sense but not in any operational one.

On every family, S-PR does worse than the layer it borrows from. Restarting the transfer walk from the allowance walk drops the transfer coefficient to $0.494$ ($\Delta\tau=-0.040$ against C-PR ($\lambda=0$)), raises the allowance coefficient only as far as $0.116$ ($\Delta\tau=-0.239$ against C-PR ($\lambda=1$)), and pushes Sybil-adjusted stability below both layers. A restart vector changes where the walk lands after each teleport, but on a graph of this density it does not change where the walk spends its time. This is the same-window counterpart of the finding that Section~\ref{sec:holdout-results} reaches out of window.

As fixed on 14~September, the secondary criterion asks C-PR's point estimate to lie within the 95\% interval of the better layer on the transfer and allowance families, and C-PR to exceed both layers on Sybil-adjusted stability. At $\lambda=0.5$, C-PR passes two of the three parts and fails the third. Its transfer coefficient ($0.527$) falls inside the interval of C-PR ($\lambda=0$) ($[0.522,0.547]$), and its Sybil-adjusted stability exceeds both layers. The allowance coefficient ($-0.001$), however, lies far outside the interval of C-PR ($\lambda=1$) ($[0.322,0.384]$), so the secondary rule fails. The paired contrasts of Table~\ref{tab:tau-diff-hybrid}, reported alongside the rule as a stricter check, would fail the transfer part as well, since they place C-PR slightly but significantly below its transfer layer. In the same-window tables the denser layer dominates the mixture. The holdout takes up the remaining question of whether the mixture has any predictive content that the single layers lack.

\dissertationSubheading[sec:holdout-results]{Temporal holdout (RQ4, RQ5)}

A same-window check cannot separate a method from its own construct. The holdout instead tests whether centrality measured up to a cut-off predicts behavior after it. Following Section~\ref{sec:holdout-protocol}, scores on the spender cohort are frozen at 28~February~2026, and two labels are then counted for each wallet between 1~March and 31~May~2026: new owner--spender approval pairs (the endorsement construct) and new transfer senders (the flow construct). Table~\ref{tab:holdout-spenders} sets every score computed at $t_1$, the two raw-degree baselines included, against both labels. Table~\ref{tab:holdout-diff} gives the paired contrasts, first those fixed before the coupled scores were computed and then those computed after the labels were known.

% Real local parquet -- regenerate with:
%   2-Dissertation-Draft-Works/wallet-reputation-experiments/scripts/run_holdout_eval.py --cohort spenders
% Generated by export_latex_results.py -- do not edit by hand unless re-export fails
\begin{table}[htbp]
\centering
\caption[Temporal holdout on the spender cohort ($n=1{,}335$): Kendall $\tau$ between scores frozen at 28~February~2026 and two labels counted on the wallet between 1~March~2026 and 31~May~2026: the number of new owner--spender approval pairs and the number of addresses that sent a transfer to the wallet for the first time.]{Temporal holdout on the spender cohort ($n=1{,}335$): Kendall $\tau$ between scores frozen at 28~February~2026 and two labels counted on the wallet between 1~March~2026 and 31~May~2026: the number of new owner--spender approval pairs and the number of addresses that sent a transfer to the wallet for the first time. CI = 95\% percentile bootstrap over wallets (400 paired resamples). 222 wallets received at least one new pair and 175 at least one new sender. Baseline rows are the raw degrees at $t_1$ that each single-layer PageRank smooths. C-PR ($\lambda=1$) and C-PR ($\lambda=0$) are the allowance and transfer layers of the coupled operator walked alone.}
\label{tab:holdout-spenders}
\fitwidth{%
\begin{tabular}{lcc}
\toprule
Score at $t_1$ & New approval pairs ($\tau$ [95\% CI]) & New transfer senders ($\tau$ [95\% CI]) \\
\midrule
In-approve degree at $t_1$ (baseline) & 0.711 [0.665, 0.759] & 0.404 [0.348, 0.467] \\
Transfer in-degree at $t_1$ (baseline) & 0.117 [0.049, 0.191] & 0.472 [0.423, 0.520] \\
\midrule
EndorseRank & 0.394 [0.360, 0.431] & 0.283 [0.247, 0.320] \\
AWP & 0.123 [0.068, 0.177] & 0.344 [0.307, 0.383] \\
\midrule
C-PR ($\lambda=1$) & 0.479 [0.427, 0.529] & 0.372 [0.325, 0.424] \\
C-PR ($\lambda=0$) & 0.044 [-0.011, 0.098] & 0.270 [0.235, 0.310] \\
C-PR ($\lambda=0.25$) & 0.271 [0.231, 0.314] & 0.266 [0.231, 0.306] \\
C-PR ($\lambda=0.5$) & 0.295 [0.252, 0.336] & 0.277 [0.241, 0.318] \\
C-PR ($\lambda=0.75$) & 0.314 [0.271, 0.357] & 0.287 [0.253, 0.330] \\
S-PR & 0.114 [0.052, 0.164] & 0.264 [0.229, 0.304] \\
\bottomrule
\end{tabular}
}
\end{table}


% Real local parquet -- regenerate with:
%   A-Skill-Programs/margin_rank/scripts/run_holdout_eval.py --cohort spenders
% Generated by export_latex_results.py -- do not edit by hand unless re-export fails
\begin{table}[htbp]
\centering
\caption[Paired bootstrap differences on the spender holdout.]{Paired bootstrap differences on the spender holdout ($n=1{,}335$, 400 resamples shared by $a$ and $b$). The first two rows give the increment of each PageRank over its own raw degree at $t_1$; the next two compare the single-layer methods on the other construct's label; the last four are the pre-registered hybrid contrasts. Primary criterion (C-PR no worse than EndorseRank on new approval pairs and better than AWP on new transfer senders): not met.}
\label{tab:holdout-diff}
\fitwidth{%
\begin{tabular}{llccccc}
\toprule
Contrast ($a$ minus $b$) & Label & $\tau_a$ & $\tau_b$ & $\Delta\tau$ & 95\% CI of $\Delta\tau$ & $P(\Delta\tau>0)$ \\
\midrule
EndorseRank minus t1 in-approve degree & new approval pairs & 0.479 & 0.711 & -0.233 & [-0.291, -0.183] & 0.000 \\
AWP minus t1 in-degree & new transfer senders & 0.270 & 0.472 & -0.201 & [-0.230, -0.171] & 0.000 \\
EndorseRank minus AWP & new approval pairs & 0.479 & 0.044 & +0.435 & [0.365, 0.511] & 1.000 \\
AWP minus EndorseRank & new transfer senders & 0.270 & 0.372 & -0.102 & [-0.148, -0.051] & 0.000 \\
C-PR minus EndorseRank & new approval pairs & 0.295 & 0.479 & -0.184 & [-0.213, -0.149] & 0.000 \\
C-PR minus AWP & new transfer senders & 0.277 & 0.270 & +0.007 & [-0.012, 0.028] & 0.735 \\
S-PR minus EndorseRank & new approval pairs & 0.114 & 0.479 & -0.365 & [-0.431, -0.300] & 0.000 \\
S-PR minus AWP & new transfer senders & 0.264 & 0.270 & -0.006 & [-0.017, 0.006] & 0.163 \\
\bottomrule
\end{tabular}
}
\end{table}


We take the three findings in turn.

The first finding is that each edge type predicts its own construct better than the other does. On new approval pairs EndorseRank reaches $\tau=0.394$ (interval $[0.360,0.431]$) against $0.123$ for AWP, and the paired difference is $+0.271$ ($[0.207,0.332]$). On new transfer senders the order reverses: AWP reaches $0.344$ against $0.283$ for EndorseRank, a difference of $+0.061$ ($[0.026,0.098]$). Every label falls after all scored events, so this is out-of-window evidence that endorsement centrality carries information about future endorsement that transfer centrality does not, and the reverse for future inflow. On this cohort EndorseRank leaves almost no ties: it gives $1{,}291$ distinct scores to the $1{,}335$ spenders, and its largest tie block holds $42$ of them. The coefficient therefore rests on an ordering of nearly every spender. The allowance layer walked alone with raw amounts, C-PR ($\lambda=1$), ranks both labels higher than EndorseRank ($0.479$ and $0.372$). Its rows are decided by unlimited approvals (Section~\ref{sub:coupled-graph}), so on this cohort the raw amounts kept part of the signal that counting each approval once gives up; the thesis does not test why.

Coverage explains much of the gap between the two walks and none of the gap between AWP and its degree. Of the 1,335 spenders, 318 have no transfer edge up to the freeze and score zero under AWP. They gained new approval pairs more often than the rest (22\% against 15\%) and new senders less often (4\% against 16\%), so AWP's zeros cost it on the first label and help it on the second. On the 977 spenders that had received a transfer from another address before the freeze (Table~\ref{tab:holdout-receiving}), EndorseRank's lead on new approval pairs shrinks to $+0.042$ (interval $[0.003,0.085]$), and the two walks are level on new senders ($-0.029$ for EndorseRank minus AWP, $[-0.070,0.011]$). AWP stays far below the transfer in-degree on that subset ($0.388$ against $0.701$; $-0.314$, $[-0.353,-0.276]$). The restriction was chosen after the results were known.

% Post hoc checks -- regenerate with:
%   2-Dissertation-Draft-Works/wallet-reputation-experiments/scripts/run_supplementary_checks.py
% Generated by export_latex_results.py -- do not edit by hand unless re-export fails
\begin{table}[htbp]
\centering
\caption[Spender holdout restricted to the 977 of the 1{,}335 spenders that had received a transfer from another address before the freeze.]{Spender holdout restricted to the 977 of the 1{,}335 spenders that had received a transfer from another address before the freeze. The other 358 have no inbound transfer edge, and 318 spenders score zero under AWP. The restriction was tried after the results were known. CI = 95\% percentile bootstrap over wallets (400 resamples, seed 42).}
\label{tab:holdout-receiving}
\fitwidth{%
\begin{tabular}{lcc}
\toprule
Score at $t_1$ & New approval pairs ($\tau$ [95\% CI]) & New transfer senders ($\tau$ [95\% CI]) \\
\midrule
EndorseRank & 0.340 [0.299, 0.386] & 0.358 [0.316, 0.396] \\
AWP & 0.298 [0.252, 0.343] & 0.388 [0.348, 0.430] \\
\midrule
In-approve degree at $t_1$ & 0.679 [0.616, 0.736] & 0.570 [0.505, 0.635] \\
Transfer in-degree at $t_1$ & 0.569 [0.510, 0.631] & 0.701 [0.651, 0.750] \\
\midrule
EndorseRank minus AWP (new approval pairs) & \multicolumn{2}{c}{$\Delta\tau=+0.042$ [0.003, 0.085]} \\
EndorseRank minus AWP (new transfer senders) & \multicolumn{2}{c}{$\Delta\tau=-0.029$ [-0.070, 0.011]} \\
AWP minus t1 in-degree (new transfer senders) & \multicolumn{2}{c}{$\Delta\tau=-0.314$ [-0.353, -0.276]} \\
\bottomrule
\end{tabular}
}
\end{table}


The second finding answers the second part of RQ4 in the negative: neither PageRank improves on its own raw degree at $t_1$. In-approve degree at $t_1$ predicts new approval pairs with $\tau=0.711$, against EndorseRank's $0.394$ ($\Delta\tau=-0.317$, interval $[-0.354,-0.283]$), and transfer in-degree at $t_1$ predicts new transfer senders with $\tau=0.472$, against AWP's $0.344$ ($\Delta\tau=-0.127$, interval $[-0.150,-0.099]$). Both intervals lie well below zero. Within a three-month window, a wallet gains more new counterparties the more it already has, and the local count captures this directly. PageRank replaces the count with a weighted, propagated quantity. In EndorseRank an owner's endorsement is shared among all the spenders it approves, and in AWP a sender's among all the addresses it pays; on these labels that sharing acts as noise. EndorseRank's out-of-window predictive content is contained in, and diluted from, the spender's $t_1$ in-approve degree, and AWP stands in the same relation to transfer in-degree. If either PageRank adds anything beyond its degree, it does not show in predicting the future growth of its own construct. What remains of EndorseRank's advantage over AWP in the first finding therefore belongs to the allowance edge, and much of it is coverage: AWP has nothing to say about spenders without a transfer history. The in-approve degree, which uses no transfer information at all, is also ahead of AWP on new senders ($0.404$ against $0.344$, without a paired test).

Third, the coupled operator fails the rule of 14~September, although in an exploratory comparison it improves on its transfer layer walked alone. That rule asked C-PR to match C-PR ($\lambda=1$) on new approval pairs and to beat C-PR ($\lambda=0$) on new transfer senders. At $\lambda=0.5$, C-PR does worse than the allowance walk on new approval pairs ($\Delta\tau=-0.184$, interval $[-0.213,-0.149]$) and no differently from the transfer walk on new transfer senders ($\Delta\tau=+0.007$, interval $[-0.012,0.028]$), so both conditions of its primary part fail. Against the transfer walk alone, a comparison the rule did not include, C-PR does better. On new approval pairs C-PR reaches $0.295$ against $0.044$, a paired difference of $+0.251$ (interval $[0.203,0.306]$, on the same resamples), and on new senders it shows no detectable loss. This contrast was computed after the labels were known, so it is treated as a hypothesis and tested again in the registered window below. Against EndorseRank and AWP, also computed after the labels were known, C-PR sits between the two on new approval pairs, $0.099$ below EndorseRank ($[-0.133,-0.064]$) and $0.172$ above AWP ($[0.111,0.226]$), and it falls behind AWP on new senders ($-0.067$, $[-0.095,-0.039]$). On the spring data the analogue of F2 would therefore fail against AWP. The registered replication reports that comparison next to its decision (Section~\ref{sub:fresh-holdout}). Across the three values of $\lambda$ the order is monotone on both labels. On new approval pairs $\tau$ rises from $0.271$ ($\lambda=0.25$) through $0.295$ to $0.314$ ($\lambda=0.75$), and on new transfer senders from $0.266$ through $0.277$ to $0.287$. More allowance weight helps on both labels; for this cohort the transfer layer adds no out-of-window signal to what the allowance layer carries. The seeded ablation shows where the allowance signal has to sit. S-PR matches the transfer walk on the flow label ($\Delta\tau=-0.006$, interval $[-0.017,0.006]$) and falls far below the allowance walk on the endorsement label ($\Delta\tau=-0.365$). When transfer edges decide where the walk goes, a restart distribution derived from endorsements changes very little. On the spender holdout, the allowance signal survives when allowance edges carry the walk, fully in the allowance walk and diluted in C-PR, and it is lost when they only decide where the walk restarts.

Two caveats apply to the whole table. Each label is the same construction carried into the future (new approvals, new senders), so none of these results bears on credit or loss. Both labels are also sparse. Of the $1{,}335$ wallets, $222$ gained a new approval pair and $175$ a new sender; every coefficient rests entirely on how that small subset is ordered.

Other candidate labels were computed on the holdout, and they are reported here without being offered as support. Revocation counts and the drain heuristic both rise with EndorseRank ($\tau=0.269$, interval $[0.226,0.308]$, and $0.422$, interval $[0.389,0.455]$), and more steeply still with the raw in-approve degree ($0.483$ and $0.575$). For an endorsement reading the sign is wrong, since a trusted spender should lose approvals less often, not more. The likelier reading is volume. A spender with many owners sees more of everything owners do, revocations and drains included, so these labels cannot separate trust from use. Aave liquidations are empty in this spender set. The same-window trading families of the research proposal (Section~\ref{sub:gf-pr}) are archived in Table~\ref{tab:alignment-summary}. Both scores align weakly with them, and AWP is ahead of EndorseRank on liquidation ($0.141$ against $0.093$) and trading success ($0.236$ against $0.096$). On the inverse-risk family, after the sign correction described in Appendix~\ref{sec:supplementary-summary}, both are negative: $-0.055$ for EndorseRank and $-0.127$ for AWP. The summed loss and the worst close get larger the more a wallet trades, and on the non-loss close rate, a share that does not, the two coefficients are $-0.031$ and $-0.006$.

\dissertationSubsubheading[sub:trader-results]{Trader outcomes (exploratory)}

The trader run of Section~\ref{sub:trader-holdout} points the other way from the spender results. It was scored after its labels were known and is reported as exploratory. Among the 1,860 matched wallets with at least three closes from March to May, AWP ranks the number of profitable closes at $\tau=0.210$ (interval $[0.181,0.238]$) against $0.072$ ($[0.041,0.111]$) for EndorseRank, and the realized gain at $0.191$ ($[0.159,0.219]$) against $0.048$ ($[0.017,0.082]$). On the share of closes that were profitable neither score shows anything: $0.016$ ($[-0.015,0.044]$) for AWP and $-0.016$ ($[-0.053,0.024]$) for EndorseRank. AWP's lead is on counts and sums, which grow with how much a wallet trades, and not on how well it trades; the transfer in-degree does about as well on the same two labels ($0.198$ and $0.169$). C-PR lies between the two scores ($0.121$ and $0.151$). The Aave liquidation label of the same run covers only 174 wallets, 10 of them liquidated, which is too few to read. That is what a transfer score should do, and it is also what one expects from a spender-side score that leaves most traders tied. The registered window gives C-PR a trader label of its own.

\dissertationSubsubheading[sub:fresh-results]{Registered replication: June to August 2026}

The registration of Section~\ref{sub:fresh-holdout} was committed to the public repository on 1~October~2026, and the June--August logs were extracted later that day. Nine monthly queries over approvals, transfers and GMX events scanned about 1.04\,TB and returned 163,553 approval logs, 1,013,023 transfer logs and 144,152 GMX closes, all of which decoded. The extraction manifest records the hashes of both registration files and no use of the testing override. Before any fresh number was computed, the evaluation code was run on the spring window as a check. It reproduced exactly the sixteen spring coefficients of the eight scores it shares with Table~\ref{tab:holdout-spenders} (the two layers walked alone, C-PR at three values of $\lambda$, S-PR and the two degrees) and the contrasts of the 14~September rule, and it gave the difference of $+0.251$ quoted above for C-PR against its transfer layer on new approval pairs. The registered evaluation then ran with 400 paired resamples and seed 42.

At the freeze on 31~May, 1,964 spenders held a positive latest allowance; 158 of them gained a new approval pair from June to August and 136 a new sender. Of the 5,521 matched wallets, 5,151 were in the freeze-date graph, and 1,402 of these closed at least three positions in the label window. Of the 1,402, 789 had at least one liquidation and 17 had every close liquidated. Table~\ref{tab:fresh-holdout} gives every registered score on the three registered labels, and Table~\ref{tab:fresh-contrasts} gives the contrasts with their verdicts. The registration documents call the two layers walked alone EndorseRank and AWP; the tables use the names of this thesis.

% Registered replication, June-August 2026 -- regenerate with:
%   2-Dissertation-Draft-Works/wallet-reputation-experiments/scripts/run_fresh_holdout.py
% Generated by export_latex_results.py -- do not edit by hand unless re-export fails
\begin{table}[htbp]
\centering
\caption[Registered replication: Kendall $\tau$ between scores frozen at 31~May~2026 and the three registered labels, counted from 1~June to 31~August~2026.]{Registered replication: Kendall $\tau$ between scores frozen at 31~May~2026 and the three registered labels, counted from 1~June to 31~August~2026. The spenders are every spender holding a positive latest allowance at the freeze; 158 of them gained a new approval pair and 136 a new sender. The traders are the 5{,}151 matched wallets in the freeze-date graph; the label is defined for the 1{,}402 with at least three closes in the label window, 789 of which had at least one liquidation. The published form of AWP belongs to a registered sensitivity analysis. CI = 95\% percentile bootstrap over wallets (400 paired resamples, seed 42).}
\label{tab:fresh-holdout}
\fitwidth{%
\begin{tabular}{lccc}
\toprule
 & \multicolumn{2}{c}{Spenders ($n=1{,}964$)} & Traders ($n=1{,}402$) \\
\cmidrule(lr){2-3} \cmidrule(lr){4-4}
Score at 31 May 2026 & New approval pairs & New transfer senders & Liquidation-free close rate \\
\midrule
In-approve degree at the freeze (baseline) & 0.613 [0.564, 0.663] & 0.371 [0.313, 0.429] & 0.144 [0.110, 0.172] \\
Transfer in-degree at the freeze (baseline) & 0.128 [0.064, 0.188] & 0.435 [0.401, 0.468] & -0.006 [-0.046, 0.031] \\
\midrule
EndorseRank & 0.338 [0.292, 0.388] & 0.289 [0.247, 0.332] & 0.086 [0.037, 0.129] \\
AWP & 0.058 [0.012, 0.103] & 0.257 [0.229, 0.286] & 0.065 [0.030, 0.102] \\
AWP, published form & 0.118 [0.071, 0.161] & 0.319 [0.291, 0.345] & 0.045 [0.009, 0.079] \\
\midrule
C-PR ($\lambda=0.25$) & 0.228 [0.190, 0.265] & 0.245 [0.217, 0.273] & 0.077 [0.042, 0.114] \\
C-PR ($\lambda=0.5$) & 0.238 [0.200, 0.274] & 0.246 [0.218, 0.274] & 0.083 [0.047, 0.119] \\
C-PR ($\lambda=0.75$) & 0.243 [0.206, 0.280] & 0.248 [0.219, 0.275] & 0.086 [0.051, 0.121] \\
S-PR & 0.126 [0.076, 0.172] & 0.254 [0.223, 0.281] & 0.030 [-0.006, 0.066] \\
\bottomrule
\end{tabular}
}
\end{table}


% Registered replication, June-August 2026 -- regenerate with:
%   A-Skill-Programs/margin_rank/scripts/run_fresh_holdout.py
% Generated by export_latex_results.py -- do not edit by hand unless re-export fails
\begin{table}[htbp]
\centering
\caption[Registered replication: the paired contrasts F1--F6 with their verdicts, and the two sensitivity analyses added before the extraction.]{Registered replication: the paired contrasts F1--F6 with their verdicts, and the two sensitivity analyses added before the extraction. $\Delta\tau=\tau_a-\tau_b$, with C-PR as $a$. Superiority holds when the interval lies above zero; non-inferiority holds when its lower end lies above $-0.02$; F6a, the first part of the 14~September rule, holds when the interval contains zero or lies above it. F1--F3 are on spenders (new approval pairs, new transfer senders) and traders (liquidation-free close rate) as in Table~\ref{tab:fresh-holdout}. CI = 95\% percentile bootstrap over wallets (400 paired resamples, seed 42).}
\label{tab:fresh-contrasts}
\fitwidth{%
\begin{tabular}{lllcccll}
\toprule
 & Contrast & Label & Test & $\tau_a$ & $\tau_b$ & $\Delta\tau$ [95\% CI] & Verdict \\
\midrule
\multicolumn{8}{l}{\emph{Registered rule: C-PR ($\lambda=0.5$) against the AWP of this thesis}} \\
F1 & C-PR vs.\ AWP & new approval pairs & superiority & 0.238 & 0.058 & +0.179 [0.132, 0.225] & holds \\
F2 & C-PR vs.\ AWP & new transfer senders & non-inferiority, $-0.02$ & 0.246 & 0.257 & -0.011 [-0.025, 0.003] & fails \\
F3 & C-PR vs.\ AWP & liquidation-free close rate & non-inferiority, $-0.02$ & 0.083 & 0.065 & +0.018 [0.011, 0.024] & holds \\
\multicolumn{7}{l}{Decision: F1, F2 and F3 must all hold} & not met \\
\midrule
\multicolumn{8}{l}{\emph{Reported, not part of the decision}} \\
F4 & C-PR vs.\ AWP & liquidation-free close rate & superiority & 0.083 & 0.065 & +0.018 [0.011, 0.024] & holds \\
F5 & C-PR vs.\ EndorseRank & liquidation-free close rate & superiority & 0.083 & 0.086 & -0.003 [-0.060, 0.051] & fails \\
F6a & C-PR vs.\ EndorseRank & new approval pairs & no worse & 0.238 & 0.338 & -0.101 [-0.127, -0.076] & fails \\
F6b & C-PR vs.\ AWP & new transfer senders & superiority & 0.246 & 0.257 & -0.011 [-0.025, 0.003] & fails \\
\midrule
\multicolumn{8}{l}{\emph{Sensitivity analysis: AWP in its published form as the comparator}} \\
F1 & C-PR vs.\ AWP, published & new approval pairs & superiority & 0.238 & 0.118 & +0.119 [0.075, 0.169] & holds \\
F2 & C-PR vs.\ AWP, published & new transfer senders & non-inferiority, $-0.02$ & 0.246 & 0.319 & -0.073 [-0.092, -0.055] & fails \\
F3 & C-PR vs.\ AWP, published & liquidation-free close rate & non-inferiority, $-0.02$ & 0.083 & 0.045 & +0.037 [0.011, 0.066] & holds \\
F4 & C-PR vs.\ AWP, published & liquidation-free close rate & superiority & 0.083 & 0.045 & +0.037 [0.011, 0.066] & holds \\
\midrule
\multicolumn{8}{l}{\emph{Sensitivity analysis: wallets outside a graph kept as isolated nodes}} \\
F1 & C-PR vs.\ AWP & new approval pairs & superiority & 0.238 & 0.053 & +0.185 [0.137, 0.232] & holds \\
F2 & C-PR vs.\ AWP & new transfer senders & non-inferiority, $-0.02$ & 0.246 & 0.258 & -0.012 [-0.027, 0.003] & fails \\
F3 & C-PR vs.\ AWP & liquidation-free close rate & non-inferiority, $-0.02$ & 0.083 & 0.065 & +0.018 [0.010, 0.024] & holds \\
F4 & C-PR vs.\ AWP & liquidation-free close rate & superiority & 0.083 & 0.065 & +0.018 [0.010, 0.024] & holds \\
F5 & C-PR vs.\ EndorseRank & liquidation-free close rate & superiority & 0.083 & 0.155 & -0.072 [-0.118, -0.028] & fails \\
F6a & C-PR vs.\ EndorseRank & new approval pairs & no worse & 0.238 & 0.338 & -0.101 [-0.127, -0.076] & fails \\
F6b & C-PR vs.\ AWP & new transfer senders & superiority & 0.246 & 0.258 & -0.012 [-0.027, 0.003] & fails \\
\bottomrule
\end{tabular}
}
\end{table}


The rule is not met. F1 holds: C-PR ranks the spenders that later gained new approval pairs at $\tau=0.238$ against $0.058$ for its transfer layer walked alone, a difference of $+0.179$ (interval $[0.132,0.225]$), so the spring gain on authorization replicates. F3 holds as well. On the liquidation-free close rate of the traders C-PR reaches $0.083$ against $0.065$ for the transfer walk ($+0.018$, $[0.011,0.024]$). F2 fails. On new transfer senders C-PR reaches $0.246$ against $0.257$; the difference is $-0.011$, and the lower end of its interval, $-0.025$, lies below the margin of $-0.02$. The interval also contains zero, so the data do not show that C-PR is worse than the transfer walk on new senders. Nor do they show that it is within $0.02$ of it, which is what the rule asked. The rule needed all three hypotheses, so, as the plan said for this case, the thesis makes no claim that C-PR improves on a walk over the transfers alone.

The plan warned that F3 could pass with both coefficients close to zero. Both are small, and both intervals lie above zero (C-PR $[0.047,0.119]$, the transfer walk $[0.030,0.102]$), so each walk carries a little information about which traders avoid liquidation. The in-approve degree at the freeze ranks this label better than any walk ($0.144$, $[0.110,0.172]$); the allowance walk reaches $0.086$ and AWP $0.045$.

The results reported beside the decision follow the same order. F4 holds: on the liquidation-free close rate C-PR is better than the transfer walk, and not only no worse. F5 fails, since C-PR and the allowance walk cannot be told apart on that label ($-0.003$, $[-0.060,0.051]$). F6, the primary part of the 14~September rule, fails again on both parts: C-PR is below the allowance walk on new approval pairs ($-0.101$, $[-0.127,-0.076]$) and not above the transfer walk on new senders. C-PR's coefficients rise with $\lambda$ on both spender labels, as in the spring. S-PR matches the transfer walk on new senders ($-0.003$, $[-0.017,0.010]$) and falls far below the allowance walk on new approval pairs ($-0.213$).

Neither registered sensitivity analysis changes the verdict (Table~\ref{tab:fresh-contrasts}). With AWP as the comparator, F1 and F3 still hold ($+0.119$ and $+0.037$) and F2 fails clearly. AWP ranks new senders at $0.319$, and C-PR's difference to it is $-0.073$ ($[-0.092,-0.055]$), an interval that lies entirely below the margin. Against AWP, then, C-PR gives up a measurable part of the ranking of later inflow in exchange for its gain on later approvals. With the wallets outside a graph kept as isolated nodes, F1 to F4 keep their verdicts. F5 changes from a null to a loss. The rule moves the matched traders outside the allowance graph from the bottom of the allowance walk's order into the block of wallets without in-edges (Section~\ref{sec:rank-divergence}); the allowance walk's coefficient on the liquidation-free close rate then rises from $0.086$ to $0.155$, and C-PR falls below it ($-0.072$, $[-0.118,-0.028]$). F6 fails under this rule too.

% Registered replication, June-August 2026 -- regenerate with:
%   2-Dissertation-Draft-Works/wallet-reputation-experiments/scripts/run_fresh_posthoc.py
% Generated by export_latex_results.py -- do not edit by hand unless re-export fails
\begin{table}[htbp]
\centering
\caption[EndorseRank on the cohorts and labels of the registered replication, scored after that replication was evaluated; none of these rows enters its decision.]{EndorseRank on the cohorts and labels of the registered replication, scored after that replication was evaluated; none of these rows enters its decision. The AWP, C-PR and degree rows repeat Table~\ref{tab:fresh-holdout}, and the contrasts use the same wallet resamples; the differences between EndorseRank and AWP are in Table~\ref{tab:endorserank-awp-holdout}. CI = 95\% percentile bootstrap over wallets (400 paired resamples, seed 42).}
\label{tab:fresh-posthoc}
\fitwidth{%
\begin{tabular}{lccc}
\toprule
 & \multicolumn{2}{c}{Spenders ($n=1{,}964$)} & Traders \\
\cmidrule(lr){2-3} \cmidrule(lr){4-4}
Score at 31 May 2026 & New approval pairs & New transfer senders & Liquidation-free close rate \\
\midrule
In-approve degree at the freeze (baseline) & 0.613 [0.564, 0.663] & 0.371 [0.313, 0.429] & 0.144 [0.110, 0.172] \\
Transfer in-degree at the freeze (baseline) & 0.128 [0.064, 0.188] & 0.435 [0.401, 0.468] & -0.006 [-0.046, 0.031] \\
\midrule
EndorseRank & 0.333 [0.306, 0.364] & 0.255 [0.227, 0.284] & 0.081 [0.032, 0.123] \\
EndorseRank, restarts by approving activity & 0.325 [0.295, 0.358] & 0.284 [0.252, 0.312] & 0.077 [0.037, 0.116] \\
AWP & 0.118 [0.071, 0.161] & 0.319 [0.291, 0.345] & 0.045 [0.009, 0.079] \\
C-PR ($\lambda=0.5$) & 0.238 [0.200, 0.274] & 0.246 [0.218, 0.274] & 0.083 [0.047, 0.119] \\
\midrule
EndorseRank minus t1 in-approve degree, new approval pairs & \multicolumn{3}{l}{$\Delta\tau=-0.280$ [-0.313, -0.247]} \\
AWP minus t1 in-degree, new transfer senders & \multicolumn{3}{l}{$\Delta\tau=-0.116$ [-0.133, -0.095]} \\
C-PR minus EndorseRank, new approval pairs & \multicolumn{3}{l}{$\Delta\tau=-0.096$ [-0.128, -0.065]} \\
C-PR minus EndorseRank, liquidation-free close rate & \multicolumn{3}{l}{$\Delta\tau=+0.002$ [-0.053, 0.056]} \\
\bottomrule
\end{tabular}
}
\end{table}


EndorseRank was not among the registered scores, so Table~\ref{tab:fresh-posthoc} scores it on the same cohorts, labels and wallet resamples after the registered evaluation; none of its rows enters the decision. On the spenders EndorseRank repeats its spring pattern. It ranks new approval pairs at $0.333$ against $0.118$ for AWP ($+0.215$, interval $[0.171,0.270]$) and new senders at $0.255$ against $0.319$ for AWP (AWP minus EndorseRank $+0.064$, $[0.036,0.088]$). The in-approve degree again beats it on new approval pairs ($0.613$; EndorseRank minus the degree $-0.280$, $[-0.313,-0.247]$), AWP again falls below the transfer in-degree on new senders ($0.319$ against $0.435$; $-0.116$, $[-0.133,-0.095]$), and C-PR again sits below it on that label ($-0.096$, $[-0.128,-0.065]$). On the liquidation-free close rate of the traders EndorseRank reaches $0.081$, close to C-PR ($0.083$) and above AWP ($0.045$), but neither difference can be told from zero ($+0.035$, $[-0.017,0.087]$ against AWP; C-PR minus EndorseRank $+0.002$, $[-0.053,0.056]$). With AWP's restarts, EndorseRank gives $0.325$ and $0.284$ on the two spender labels and $0.077$ on the traders.

% Registered replication, June-August 2026 -- regenerate with:
%   2-Dissertation-Draft-Works/wallet-reputation-experiments/scripts/run_fresh_holdout.py
% Generated by export_latex_results.py -- do not edit by hand unless re-export fails
\begin{table}[htbp]
\centering
\caption[Registered replication, traders: Kendall $\tau$ between scores frozen at 31~May~2026 and the other GMX~V2 labels counted from June to August~2026, on the 1{,}402 matched wallets with at least three closes in that window.]{Registered replication, traders: Kendall $\tau$ between scores frozen at 31~May~2026 and the other GMX~V2 labels counted from June to August~2026, on the 1{,}402 matched wallets with at least three closes in that window. No liquidation is 1 when none of the wallet's closes was a liquidation. Profitable closes and realized gain grow with the number of closes; the two shares do not. These labels are reported and do not enter the decision. CI = 95\% percentile bootstrap over wallets (400 paired resamples, seed 42).}
\label{tab:fresh-traders}
\fitwidth{%
\begin{tabular}{lccccc}
\toprule
Score at 31 May 2026 & No liquidation & Profitable closes & Realized gain & Profitable share & Non-loss share \\
\midrule
In-approve degree at the freeze (baseline) & \begin{tabular}[c]{@{}c@{}}0.187\\[-1pt]{\footnotesize [0.150, 0.221]}\end{tabular} & \begin{tabular}[c]{@{}c@{}}-0.004\\[-1pt]{\footnotesize [-0.033, 0.028]}\end{tabular} & \begin{tabular}[c]{@{}c@{}}0.021\\[-1pt]{\footnotesize [-0.015, 0.054]}\end{tabular} & \begin{tabular}[c]{@{}c@{}}-0.074\\[-1pt]{\footnotesize [-0.105, -0.042]}\end{tabular} & \begin{tabular}[c]{@{}c@{}}-0.139\\[-1pt]{\footnotesize [-0.172, -0.101]}\end{tabular} \\
Transfer in-degree at the freeze (baseline) & \begin{tabular}[c]{@{}c@{}}-0.072\\[-1pt]{\footnotesize [-0.115, -0.030]}\end{tabular} & \begin{tabular}[c]{@{}c@{}}0.191\\[-1pt]{\footnotesize [0.150, 0.224]}\end{tabular} & \begin{tabular}[c]{@{}c@{}}0.127\\[-1pt]{\footnotesize [0.089, 0.166]}\end{tabular} & \begin{tabular}[c]{@{}c@{}}0.026\\[-1pt]{\footnotesize [-0.008, 0.060]}\end{tabular} & \begin{tabular}[c]{@{}c@{}}-0.001\\[-1pt]{\footnotesize [-0.036, 0.038]}\end{tabular} \\
\midrule
EndorseRank & \begin{tabular}[c]{@{}c@{}}0.069\\[-1pt]{\footnotesize [0.011, 0.116]}\end{tabular} & \begin{tabular}[c]{@{}c@{}}0.074\\[-1pt]{\footnotesize [0.036, 0.107]}\end{tabular} & \begin{tabular}[c]{@{}c@{}}0.016\\[-1pt]{\footnotesize [-0.024, 0.059]}\end{tabular} & \begin{tabular}[c]{@{}c@{}}-0.043\\[-1pt]{\footnotesize [-0.084, -0.000]}\end{tabular} & \begin{tabular}[c]{@{}c@{}}-0.130\\[-1pt]{\footnotesize [-0.166, -0.091]}\end{tabular} \\
AWP & \begin{tabular}[c]{@{}c@{}}0.014\\[-1pt]{\footnotesize [-0.026, 0.056]}\end{tabular} & \begin{tabular}[c]{@{}c@{}}0.140\\[-1pt]{\footnotesize [0.103, 0.177]}\end{tabular} & \begin{tabular}[c]{@{}c@{}}0.145\\[-1pt]{\footnotesize [0.110, 0.180]}\end{tabular} & \begin{tabular}[c]{@{}c@{}}0.033\\[-1pt]{\footnotesize [-0.001, 0.069]}\end{tabular} & \begin{tabular}[c]{@{}c@{}}-0.004\\[-1pt]{\footnotesize [-0.040, 0.035]}\end{tabular} \\
AWP, published form & \begin{tabular}[c]{@{}c@{}}-0.030\\[-1pt]{\footnotesize [-0.069, 0.011]}\end{tabular} & \begin{tabular}[c]{@{}c@{}}0.233\\[-1pt]{\footnotesize [0.193, 0.268]}\end{tabular} & \begin{tabular}[c]{@{}c@{}}0.182\\[-1pt]{\footnotesize [0.142, 0.217]}\end{tabular} & \begin{tabular}[c]{@{}c@{}}0.030\\[-1pt]{\footnotesize [-0.004, 0.060]}\end{tabular} & \begin{tabular}[c]{@{}c@{}}-0.035\\[-1pt]{\footnotesize [-0.065, -0.001]}\end{tabular} \\
\midrule
C-PR ($\lambda=0.25$) & \begin{tabular}[c]{@{}c@{}}0.032\\[-1pt]{\footnotesize [-0.007, 0.072]}\end{tabular} & \begin{tabular}[c]{@{}c@{}}0.133\\[-1pt]{\footnotesize [0.094, 0.168]}\end{tabular} & \begin{tabular}[c]{@{}c@{}}0.145\\[-1pt]{\footnotesize [0.110, 0.182]}\end{tabular} & \begin{tabular}[c]{@{}c@{}}0.027\\[-1pt]{\footnotesize [-0.007, 0.062]}\end{tabular} & \begin{tabular}[c]{@{}c@{}}-0.011\\[-1pt]{\footnotesize [-0.047, 0.028]}\end{tabular} \\
C-PR ($\lambda=0.5$) & \begin{tabular}[c]{@{}c@{}}0.043\\[-1pt]{\footnotesize [0.004, 0.082]}\end{tabular} & \begin{tabular}[c]{@{}c@{}}0.124\\[-1pt]{\footnotesize [0.084, 0.160]}\end{tabular} & \begin{tabular}[c]{@{}c@{}}0.142\\[-1pt]{\footnotesize [0.107, 0.179]}\end{tabular} & \begin{tabular}[c]{@{}c@{}}0.024\\[-1pt]{\footnotesize [-0.010, 0.058]}\end{tabular} & \begin{tabular}[c]{@{}c@{}}-0.010\\[-1pt]{\footnotesize [-0.045, 0.028]}\end{tabular} \\
C-PR ($\lambda=0.75$) & \begin{tabular}[c]{@{}c@{}}0.053\\[-1pt]{\footnotesize [0.011, 0.091]}\end{tabular} & \begin{tabular}[c]{@{}c@{}}0.110\\[-1pt]{\footnotesize [0.068, 0.147]}\end{tabular} & \begin{tabular}[c]{@{}c@{}}0.135\\[-1pt]{\footnotesize [0.101, 0.171]}\end{tabular} & \begin{tabular}[c]{@{}c@{}}0.021\\[-1pt]{\footnotesize [-0.014, 0.056]}\end{tabular} & \begin{tabular}[c]{@{}c@{}}-0.008\\[-1pt]{\footnotesize [-0.041, 0.030]}\end{tabular} \\
S-PR & \begin{tabular}[c]{@{}c@{}}-0.030\\[-1pt]{\footnotesize [-0.070, 0.014]}\end{tabular} & \begin{tabular}[c]{@{}c@{}}0.176\\[-1pt]{\footnotesize [0.139, 0.211]}\end{tabular} & \begin{tabular}[c]{@{}c@{}}0.132\\[-1pt]{\footnotesize [0.094, 0.170]}\end{tabular} & \begin{tabular}[c]{@{}c@{}}0.040\\[-1pt]{\footnotesize [0.003, 0.075]}\end{tabular} & \begin{tabular}[c]{@{}c@{}}-0.023\\[-1pt]{\footnotesize [-0.057, 0.015]}\end{tabular} \\
\bottomrule
\end{tabular}
}
\end{table}


Table~\ref{tab:fresh-traders} gives the registered trader labels that do not enter the decision, and they repeat the spring trader run. AWP ranks the number of profitable closes ($0.233$) and the realized gain ($0.182$) highest of the registered scores; the transfer in-degree reaches $0.191$ and $0.127$, and C-PR $0.124$ and $0.142$. Every score stays at or below $0.04$ on the share of closes that were profitable and at or below zero on the share that were not losses.

\dissertationSubheading[sec:empirical-claims]{Primary claims (C1--C5)}

The chapter's results support the main claims below, cross-referenced to the hypothesis and research question columns of Table~\ref{tab:hrq-claim-map}.

\begin{enumerate}
\item[C1.] On identical wallet sets EndorseRank's solver call, matrix assembly included, takes a fraction of AWP's wall-clock time and memory, because its latest-allowance graph has roughly one-ninth as many edges as the transfer graph and both assembly and iteration scale with $|E|$; it also needs fewer iterations (Table~\ref{tab:benchmark-runtime}). The ratio is $9.2\times$ on the matched cohort and stays between $8.6\times$ and $14.4\times$ across in-cohort and pool stages whose edge counts differ by a factor of five to seven (Tables~\ref{tab:benchmark-scaling} and~\ref{tab:benchmark-scaling-incohort}), although the pool stages do not test edge counts beyond those of the cohort graph. The claim covers EndorseRank alone. Running on both edge sets, the coupled operator costs at least as much as AWP.
\item[C2.] On same-window local statistics each method recovers the construct it targets, and the bootstrap confidence intervals exclude zero (Table~\ref{tab:alignment-family-ci}). These tests measure internal consistency; external validity is outside their scope.
\item[C3.] Which method performs better depends on the construct. AWP is better than EndorseRank on the transfer family and EndorseRank is better than AWP on the allowance family, with paired-difference confidence intervals that exclude zero in opposite directions (Table~\ref{tab:tau-diff}), and the two score vectors agree only in part (Table~\ref{tab:alignment-family-ci}). Within EndorseRank, allowance alignment is not significantly higher than transfer alignment, but the two coefficients have different ceilings (Section~\ref{sec:rank-divergence}), so this null says nothing either way. With the wallets outside the allowance graph kept as isolated nodes, the difference is large and positive (Table~\ref{tab:tie-sensitivity}).
\item[C4.] Out of window, each edge type carries information about its own construct that the other lacks. On the spring spender holdout EndorseRank beats AWP on new approvals ($+0.271$) and AWP beats EndorseRank on new senders ($+0.061$) (Table~\ref{tab:holdout-diff}). In the June--August window, where EndorseRank was scored after the registered evaluation, the same pattern holds: EndorseRank leads on new approvals ($+0.215$) and AWP on new senders ($+0.064$) (Table~\ref{tab:fresh-posthoc}). The gain lies in the edge and not in propagation, because the spender's raw in-approve degree at $t_1$ predicts new approvals better than EndorseRank does, and the transfer in-degree predicts new senders better than AWP, in both windows. On the matched traders, AWP is ahead on the same-window transfer and stability proxies and, in the exploratory spring run, on counts of later profitable closes. None of these comparisons was fixed in advance.
\item[C5.] A single random walk over both edge types does not beat its two layers walked alone, and the registered replication does not confirm that it improves on the transfer walk. On new approval pairs C-PR falls between the two walks at every $\lambda$ tested, and the rule of 14~September is not met; the seeded variant reproduces the transfer walk (Tables~\ref{tab:holdout-spenders} and~\ref{tab:holdout-diff}). Against the transfer walk alone, C-PR at $\lambda=0.5$ adds $0.251$ on new approval pairs in the spring with no detectable loss on new senders, a comparison made after the fact. In the registered June--August window C-PR again beats the transfer walk on new approval pairs (F1) and ranks liquidation-free traders slightly better (F3, F4), but its difference on new senders cannot be bounded within $0.02$ (F2), so the registered rule is not met (Table~\ref{tab:fresh-contrasts}). Against AWP the result is clearer still: in the registered sensitivity analysis C-PR loses $0.073$ on new senders, and in the spring it sits between EndorseRank and AWP on new approvals and below AWP on new senders. Section~\ref{sec:coupled-results} reports where the hybrids stand within the window.
\end{enumerate}

The inter-method agreement and Figure~\ref{fig:rank-scatter} answer RQ1, the family tables and difference tests RQ2, the benchmark tables RQ3, and the holdouts RQ4 and RQ5. None of these claims retires AWP, and none promotes EndorseRank past what the intervals allow. EndorseRank is cheaper and aligns with the allowance checks, and its moderate agreement with AWP comes from the wallets it scores zero. Out of window it beats AWP on new approvals in both windows and trails it on new senders, yet the raw degree it smooths does better still. Mixing the two layers in one walk kept part of the allowance signal on new approvals but gave up ground on new senders, so the thesis makes no claim that the coupled walk improves on a walk over the transfers alone, or on AWP.

Chapter~\ref{ch:discussion} reads these findings against the research questions and objectives set in Chapter~\ref{ch:introduction}. It draws out their implications for theory, practice and society and then turns to the threats to validity and the limitations that bound them.
